%% =====================================================================
%%  Where that exposition departs from arXiv:2404.01080v2, and why.
%%
%%  GENERATED FILE -- do not edit.  Produced by make_standalone.py from
%%  csp-audit.tex and the per-section sources; edit those and regenerate.
%%  Line references in a review should therefore be resolved back to the
%%  section sources, whose names appear in the banners below.
%%  Generated from revision b919510 on 2026-08-02.
%% =====================================================================

%% ---------------------------------------------------------------- csp-audit.tex
\documentclass[11pt]{article}
\def\hazardheading{Note}
\def\docheading{CSP dichotomy --- source audit}
%% ---------------------------------------------------------------- preamble.tex
% Shared preamble for the three CSP dichotomy documents: the proof paper
% (csp-proof.tex), the source audit (csp-audit.tex), and the formalization
% blueprint (csp-blueprint.tex).
%
% House style follows the earlier `zhuk_centers` blueprint: numbered theorem
% environments shared across a single counter, a `longtheorem` upright variant
% for statements with many numbered hypotheses, and explicit `convention`
% environments for standing hypotheses.
%
% Each main file may set \hazardheading and \docheading before \input-ing this
% file; the defaults below apply otherwise.
\usepackage[a4paper,margin=27mm,headheight=14pt]{geometry}
\usepackage[T1]{fontenc}
\usepackage{lmodern}
\usepackage{amsmath,amssymb,amsthm}
\usepackage{longtable,array}
\usepackage{fancyhdr}
\usepackage[hidelinks]{hyperref}

\newtheoremstyle{mainplain}{9pt}{9pt}{\itshape}{}{\bfseries}{.}{0.55em}{}
\newtheoremstyle{mainroman}{9pt}{9pt}{\normalfont}{}{\bfseries}{.}{0.55em}{}
\theoremstyle{mainplain}
\newtheorem{theorem}{Theorem}[section]
\newtheorem{lemma}[theorem]{Lemma}
\newtheorem{proposition}[theorem]{Proposition}
\newtheorem{corollary}[theorem]{Corollary}
% A `claim' is a statement local to a proof, numbered in the same sequence so
% that it can be referred to from the surrounding text.
\newtheorem{claim}[theorem]{Claim}
\theoremstyle{mainroman}
\newtheorem{longtheorem}[theorem]{Theorem}
\newtheorem{longlemma}[theorem]{Lemma}
\newtheorem{longproposition}[theorem]{Proposition}
\newtheorem{longcorollary}[theorem]{Corollary}
\theoremstyle{definition}
\newtheorem{definition}[theorem]{Definition}
\newtheorem{remark}[theorem]{Remark}
\newtheorem{convention}[theorem]{Convention}
\providecommand{\hazardheading}{Formalization note}
\newtheorem{hazard}[theorem]{\hazardheading}
\newtheorem{entry}[theorem]{Entry}

\setlength{\parindent}{1.25em}\setlength{\parskip}{0pt}\setlength{\emergencystretch}{2em}
\newcommand{\alphlabels}{\renewcommand{\labelenumi}{\textup{(\alph{enumi})}}%
  \renewcommand{\theenumi}{\alph{enumi}}}
\newcommand{\romanlabels}{\renewcommand{\labelenumi}{\textup{(\roman{enumi})}}%
  \renewcommand{\theenumi}{\roman{enumi}}}
\providecommand{\tightlist}{\setlength{\itemsep}{0pt}\setlength{\parskip}{0pt}}
\numberwithin{equation}{section}
\setcounter{tocdepth}{2}\setcounter{secnumdepth}{3}
\pagestyle{fancy}\fancyhf{}
\providecommand{\docheading}{CSP dichotomy}
\fancyhead[L]{\small\docheading}\fancyhead[R]{\small\nouppercase{\leftmark}}
\fancyfoot[C]{\thepage}\renewcommand{\headrulewidth}{0.3pt}

% ------------------------------------------------------- cross-document keys
% The three documents are compiled separately, so a reference from one to a
% statement of another cannot be a \ref.  It is the statement's *label*, set in
% typewriter, which is stable and mechanically checkable.
\newcommand{\pkey}[1]{\texttt{#1}}
% A zero-width break opportunity, for the source's very long statement names.
\newcommand{\brk}{\discretionary{}{}{}}
\newcommand{\inproof}[1]{the proof paper, \pkey{#1}}

% ---------------------------------------------------------------- operators
\DeclareMathOperator{\Sg}{Sg}
\DeclareMathOperator{\Clo}{Clo}
\DeclareMathOperator{\Pol}{Pol}
\DeclareMathOperator{\Inv}{Inv}
\DeclareMathOperator{\Var}{Var}
\DeclareMathOperator{\Con}{Con}
\DeclareMathOperator{\Sol}{Sol}
\DeclareMathOperator{\Viol}{Viol}
\DeclareMathOperator{\Expanded}{Exp}
\DeclareMathOperator{\CSP}{CSP}
\DeclareMathOperator{\GenSAT}{SAT}
\DeclareMathOperator*{\proj}{pr}
\newcommand{\ConOne}{\operatorname{Con}_1}
\DeclareMathOperator{\LeftLinked}{LeftLinked}
\DeclareMathOperator{\RightLinked}{RightLinked}

% ---------------------------------------------------------------- algebras
\newcommand{\bA}{\mathbf A}
\newcommand{\bB}{\mathbf B}
\newcommand{\bC}{\mathbf C}
\newcommand{\bD}{\mathbf D}
\newcommand{\bE}{\mathbf E}
\newcommand{\bR}{\mathbf R}
\newcommand{\bS}{\mathbf S}
\newcommand{\bZ}{\mathbf Z}
\newcommand{\Vn}{\mathcal V_{n}}

% ---------------------------------------------------------- subuniverse types
% The six types of Zhuk's theory.  Typed as upright small capitals so that a
% type never reads as a product of variables.
\newcommand{\TBA}{\textup{\textsc{ba}}}
\newcommand{\TC}{\textup{\textsc{c}}}
\newcommand{\TS}{\textup{\textsc{s}}}
\newcommand{\TD}{\textup{\textsc{d}}}
\newcommand{\TL}{\textup{\textsc{l}}}
\newcommand{\TPC}{\textup{\textsc{pc}}}
\newcommand{\MT}[1]{\mathcal M#1}
\newcommand{\TML}{\MT\TL}
\newcommand{\TMPC}{\MT\TPC}
\newcommand{\TMD}{\MT\TD}

% `C <_{T}^{A} B' and its relatives.
\newcommand{\subt}[2][]{<_{#2}^{#1}}
\newcommand{\subte}[2][]{\le_{#2}^{#1}}
\newcommand{\dsubt}[2][]{\dot<_{#2}^{#1}}
\newcommand{\dsubte}[2][]{\dot\le_{#2}^{#1}}
\renewcommand{\lll}{\mathrel{\lhd\!\lhd\!\lhd}}
\newcommand{\dlll}{\mathrel{\dot{\lhd}\!\dot{\lhd}\!\dot{\lhd}}}
\newcommand{\sdle}{\le_{\mathrm{sd}}}

% congruence cover  sigma^*
\newcommand{\cover}[1]{#1^{*}}
% bridge collapse
\newcommand{\bcol}[1]{\widetilde{#1}}
% absorption
\newcommand{\ab}{\lhd}
\newcommand{\abin}{\lhd_{\mathrm{bin}}}
\newcommand{\abc}{\lhd_{\mathrm{c}}}

\hypersetup{pdftitle={Zhuk's simplified CSP dichotomy proof: a source audit},
  pdfsubject={Where our rendering departs from arXiv:2404.01080v2, and why},
  pdfauthor={}}

\title{Zhuk's Simplified CSP Dichotomy Proof\\[2pt]
  \large A source audit}
\author{}\date{}

\begin{document}
\maketitle\vspace{-2.2em}
\begin{abstract}
The companion proof paper corrects \cite{ZhukSimplified} silently: where the
source is wrong, incomplete or ambiguous, it states the corrected thing and
proves it, with nothing in the margin about what the source has. This document
is the record of those departures. For each it says what the source has, quotes
it, says why it cannot be transcribed, gives the repair, and names the
statement of the proof paper that carries it.

Sixteen substantive defects, of which fifteen are repaired with no new
mathematics and one --- the maximal multi-type extension --- is reduced to a
single stated hypothesis that remains unproved. Ten
readings legislated, two of which turn out to be theorems rather than choices.
Two claims raised and refuted, one declared moot. Four hedges in \S5 of the
source, all four honest.

References of the form \texttt{lem:\dots} are labels of statements in the proof
paper, which is compiled separately; the label is the join key between the two
documents.
\end{abstract}

\tableofcontents\clearpage

%% ---------------------------------------------------------------- audit0-intro.tex
\section{What this document is}\label{sec:audit-intro}

The proof paper is written in the Jordan--Schönflies mould: where
\cite{ZhukSimplified} is wrong, incomplete, or ambiguous, the proof paper
simply states the corrected thing and proves it, with no marginal notes about
what the source has. That is the right shape for a document whose job is to be
checked as mathematics. It is the wrong shape for a reader who wants to know
\emph{where our rendering differs from Zhuk's paper and why}, and it would be
dishonest to make those differences invisible.

This document is the difference. For every place where the proof paper departs
from \cite{ZhukSimplified}, there is an entry here that says what the source
has, quotes it, says why it cannot be transcribed, gives the repair, and names
the statement of the proof paper that carries it.

None of this is a criticism of \cite{ZhukSimplified}, which is unusually careful
prose --- across nine thousand lines of \TeX{} it says ``obvious'' three times
and ``clearly'' never, and a mechanical sweep of \S5 turns up four hedges in
thirteen printed pages, all four of which turn out to be honest
(\S\ref{sec:hedges}). Sixteen substantive defects in fifty pages is roughly
the rate one should expect of any informal exposition of a proof this size, and
finding them is the point of reading it this way.

\subsection{The evidentiary standard}\label{sec:standard}

A claim counts as checked here only when the cited lines \emph{and the prose
around them} have been read, and the passage that does or does not discharge
the claim can be quoted. The standard is not pedantry: two of the original
claims died to it. One was taken from a reader who worked from a pseudocode
block's \texttt{Input:} line while the hypothesis it needed was in the
paragraph above (\S\ref{sec:refuted-depth}); the other asserted a citation
cycle that a parser refuted in minutes (\S\ref{sec:refuted-cycle}). Both
failure modes are the same: reading the formal-looking part of a paper and not
the sentence before it.

Where a claim is reducible to a finite check, it has had one. The searches are
listed in \S\ref{sec:machine-checks} with their scope, so that a reader can see
what they do and do not cover.

\subsection{Line numbers, and how entries are keyed}\label{sec:keying}

Line numbers are into \texttt{main.tex} and \texttt{StrongSubalgebras.tex} of
the v2 arXiv source of \cite{ZhukSimplified}, obtainable from
\texttt{arxiv.org/e-print/2404.01080v2}. Zhuk's paper is not redistributed
here, and its statement names are its own: \texttt{LEMPropagation},
\texttt{THMMainStableIntersection} and so on are the source's labels, not ours.

Each entry names the statement of the proof paper that carries the repair, by
that statement's label, set in typewriter: \pkey{lem:propagation},
\pkey{thm:stable-intersection}. The two documents are compiled separately and
share no cross-reference machinery, deliberately, so that neither can be read
as evidence for the other. The label is the join key.

The tags $\mathrm D2$, $\mathrm D4,\dots$ are historical: they were assigned in
the order the claims were raised, and some early tags were retired when the
claims they named were refuted or restated. The register of
\S\ref{sec:register} is the current and complete list; the retired claims are
in \S\ref{sec:refuted}.

\subsection{Where things stand}\label{sec:audit-tally}

\begin{center}\small
\begin{tabular}{@{}p{0.52\textwidth}r@{}}
\hline
\textbf{} & \textbf{Count} \\
\hline
Defects confirmed and repaired, needing no new mathematics & 15 \\
\quad of which: a citation the later paper dropped & 3 \\
\quad a typographical or typing slip & 2 \\
\quad a statement that is false as printed & 2 \\
\quad omitted steps and dropped hypotheses & 8 \\
Defects confirmed and \textbf{open} & 1 \\
Claims examined and refuted & 2 \\
Claims examined and declared moot & 1 \\
Readings legislated (of which two are theorems, not choices) & 10 \\
Hedges in \S5 of the source, all discharged & 4 \\
\hline
\end{tabular}
\end{center}

The single open item is $\mathrm{D}17$ (\S\ref{sec:D17}), the maximal
multi-type extension. It was declared solved in an earlier draft; an external
review found the repair invalid at its first line, and the finding is correct.
It is recorded here as open rather than quietly weakened.

Its status has since been sharpened rather than resolved. The proof paper now
proves the lemma from an \emph{explicitly stated} extra hypothesis --- that
the type-$T$ factors alone already cut $B_{2}$ out --- and shows that among
those factors the image-versus-intersection failure does not occur, so the
whole difficulty is now localised in one displayed sentence rather than spread
through a proof. The hypothesis itself is unproved and is a debt at the single
call site.
\clearpage
%% ---------------------------------------------------------------- audit1-defects.tex
\section{The register}\label{sec:register}

Sixteen places where \cite{ZhukSimplified} is wrong, incomplete, or states less
than its own proof gives. Fifteen are repaired and none of those repairs needs
new mathematics; the sixteenth, $\mathrm{D}17$, is open.

\begin{center}\footnotesize
\begin{longtable}{@{}llp{0.38\textwidth}p{0.26\textwidth}@{}}
\hline
\textbf{Tag} & \textbf{Where} & \textbf{What} & \textbf{Proof paper} \\
\hline\endhead
$\mathrm D1$ & \S2.3 & Lemma~19 restates a corollary of \cite{ZhukStrong}
  dropping two hypotheses; false as printed & \pkey{lem:ba-center-implies} \\
$\mathrm D2$ & \texttt{main:1682} & Propagation modulo a congruence: stated,
  used five times, never proved & \pkey{lem:propagation} \\
$\mathrm D4$ & \texttt{SS:2594} & $n=2$ for type $\TC$ in stable intersection
  is argued only for $\TPC$ & \pkey{thm:stable-intersection} \\
$\mathrm D6$ & \S3.2 & Bridges-from-relations applied without its third
  hypothesis, which can fail & \pkey{lem:connected} \\
$\mathrm D7$ & \S3.3 & Two gaps in the main induction: \emph{connected} proved
  where \emph{linked} is demanded, and (1c) derived for the wrong reduction &
  \pkey{thm:main-inductive} \\
$\mathrm D8$ & \texttt{main:3526} & Minimal-containing-a-point is used where
  minimal-outright is required; the bridging lemma was deleted &
  \pkey{lem:minimal-consistent} \\
$\mathrm D{11}$ & \texttt{SS:1154} & The symmetrisation formula defines a
  relation that is never a bridge & \pkey{lem:symmetrise} \\
$\mathrm D{12}$ & \texttt{SS:1103} & Bridges and abelian groups: \emph{false as
  stated}, with a three-element counterexample &
  \pkey{lem:nice-bridge-abelian} \\
$\mathrm D{13}$ & \texttt{SS:1273} & Linkedness asserted from an inclusion that
  bears on a different hypothesis & \pkey{lem:nontrivial-bridge} \\
$\mathrm D{14}$ & \texttt{SS:1394} & The box form of
  absorption-preserves-linkedness cited for a restriction that is not a box &
  \pkey{lem:absorb-linked-sub} \\
$\mathrm D{15}$ & \texttt{SS:2208} & Existence of an irreducible congruence
  factors by the wrong congruence, hiding an unsupported properness step &
  \pkey{lem:main-existence} \\
$\mathrm D{16}$ & \texttt{SS:2743} & Linear congruences on top: the conclusion
  is asserted & \pkey{lem:linear-on-top} \\
$\mathrm D{17}$ & \texttt{SS:3113} & Maximal multi-type extension concludes
  that two sets meet because each meets a third. \textbf{Open} &
  \pkey{lem:maximal-mult} \\
$\mathrm C{10}$ & \texttt{SS:208} & The central-relation citation drops the
  third alternative of the theorem it cites &
  \pkey{lem:central-relation} \\
$\mathrm F$ & \texttt{SS:545, 837, 921} & A full fibre inferred from a
  surjective projection, at three sites & \pkey{lem:full-fibre} \\
$\mathrm D{18}$ & \texttt{main:1373} & The definition of a linear congruence
  is ill-typed: it intersects a relation on elements with a set of tuples of
  classes & \pkey{def:linear-congruence} \\
\hline
\end{longtable}
\end{center}

Two patterns recur rather than occurring once, and are worth naming because a
reader checking the proof paper will meet them repeatedly. The first is a
\emph{proper} containment asserted where the cited lemma gives only
$\subseteq$, usually after factoring by a congruence, which can close the gap;
it appears at least five times and is repaired the same way each time, by
producing the witness the source has already displayed two lines earlier. The
second is an upgrade from ``the image has size one'' to ``the image is
\emph{this} block''; it appears four times and the proof paper makes it a named
step.

\subsection{$\mathrm D1$ --- two hypotheses dropped, and the statement becomes false}
\label{sec:D1}

\cite[Lem.~19]{ZhukSimplified} restates \cite[Cor.~6.1.2, 6.9.2]{ZhukStrong}.
The source it cites reads ``Suppose $R\le A_{1}\times\dots\times A_{n}$,
$\proj_{1}(R)=A_{1}$, $B_{i}\le_{\TBA(t)}A_{i}$ for every $i$''. The
restatement drops the subdirectness $\proj_{1}(R)=A_{1}$, and drops --- in the
$\TBA$ case --- that the $B_{i}$ absorb with a \emph{common} term $t$, which is
the reason \cite{ZhukStrong} carries $t$ in the notation.

\emph{Without the first the statement is false.} Take
$R=\{(0,0)\}\le\mathbf Z_{p}\times\mathbf Z_{p}$, a subuniverse by
idempotence, and $C_{i}=A_{i}$, which is allowed because $\subte{T}$ admits
equality. Then $R\cap(C_{1}\times C_{2})=R$ and $\proj_{1}(R)=\{0\}$, which is
neither empty nor all of $\mathbf Z_{p}$; the conclusion therefore asserts
$\{0\}\subt{\TBA}\mathbf Z_{p}$, contradicting the paper's own Lemma~29.

\emph{Repair.} Restore both hypotheses. The proof paper states them in
\pkey{lem:ba-center-implies}. The common-term condition is not free at the
call site: \pkey{thm:main-inductive}(2) records only that every coordinate
reduction has type $\TBA$, and applies the lemma to the relation computing
$E_{2}$ from $E_{1}$. Either a lemma synthesising one common binary term from a
finite family of coordinate absorptions is needed, or the reduction data must
carry a common witness from the start. This is recorded as an obligation on the
proof paper, not as a defect of the source, since the source's own version of
the lemma carries the hypothesis.

\subsection{$\mathrm D2$ --- propagation modulo a congruence is never proved}
\label{sec:D2}

Stated at \texttt{main:1682}; used at \texttt{SS:2506, 2514, 2588, 3110} and
\texttt{main:2999}. Never proved, and not attributed to any external source.

The paper's own convention (\texttt{main:1058}) is \emph{``we always duplicate
statements if the proof appears in a later section''}, implemented by
\texttt{\textbackslash newtheorem*} duplicates. Checking all thirteen
statements of \S2.3 against that convention leaves exactly two not restated in
\S5: this corollary, and \texttt{LEMBACenterImplies}, which is declared an
external import and so is absent by design.

\emph{Repair.} It is Propagation specialised to the canonical surjection
$\pi\colon\mathbf A\twoheadrightarrow\mathbf A/\delta$, whose clauses (f),
(ft), (fs), (fm) are its four clauses. The identification is literal rather
than approximate: $X/\delta$ is \emph{defined} as the image
$\{x/\delta:x\in X\}$, which is $\pi(X)$; in particular the corollary's
hypothesis ``$B/\delta$ is $\TS$-free'' and Propagation's ``$f(B)$ is
$\TS$-free'' are the same condition. Proof paper: \pkey{cor:prop-quotient},
derived from \pkey{lem:propagation}.

A mechanical audit of all $81$ labelled statements of the source
(\S\ref{sec:machine-checks}) shows this is the \emph{only} statement that is
cited while being neither proved nor attributed. That matters beyond this
entry: the audit covers \S5.6 and \S5.7, so whatever else is wrong there, it is
not another statement asserted from nowhere.

\subsection{$\mathrm D4$ --- $n=2$ for type $\TC$ is unargued}\label{sec:D4}

The stable-intersection theorem concludes, in its case (c), that $n=2$ and
$T_{1}=T_{2}=\TC$. The proof at \texttt{SS:2594}\,ff.\ reduces general $n$ to a
pair, obtains $T_{1}=T_{2}$, and then:

\begin{quote}\small
``If $T_{1}=\TPC$ then $\sigma_{1}\circ\sigma_{2}=\sigma_{1}=\sigma_{2}$ \dots
Additionally, this implies that $n$ cannot be greater than 2 for
$T_{1}=\TPC$, as in this case the intersection $C_{1}\cap C_{2}$ must be
empty.''
\end{quote}

That argument is specific to $\TPC$: equal congruences make $C_{1}$ and $C_{2}$
distinct blocks, hence \emph{pairwise} disjoint, which contradicts hypothesis
(3) once $n\ge3$. For type $\TC$ there is no congruence and no analogue, and
the proof offers nothing; hypothesis (2) gives only
$C_{1}\cap\dots\cap C_{n}=\varnothing$, never $C_{1}\cap C_{2}=\varnothing$.
The same gap is visible at \texttt{SS:2650}, where the $\TPC$ argument needs
one further unstated step --- the two $\sigma$-blocks are distinct because
$B_{1}\cap C_{2}'\neq\varnothing$ would otherwise put a point in
$C_{1}\cap C_{2}'$ --- which again has no type-$\TC$ analogue.

\emph{Repair.} The missing argument is \cite[Thm.~3.7]{ZhukStrong}, which is
the same statement for subuniverses of a \emph{single} algebra, resting on
\cite[Lem.~6.11]{ZhukStrong}: for $n\ge3$ and $C_{i}$ central in $A_{i}$, no
$(C_{1},\dots,C_{n})$-essential relation $R\le A_{1}\times\dots\times A_{n}$
exists. The 2024 setting is not literally that of Theorem~3.7 --- there the
$B_{i}$ are subuniverses of one algebra, whereas here $C_{i}$ is central in
$B_{i}$ and the $B_{i}$ differ --- so a reduction is needed first: put
$D=\bigcap_{i}B_{i}$ and $E_{i}=C_{i}\cap D$; then $D$ is a nonempty
subuniverse, each $E_{i}$ is properly central in $D$, $\bigcap_{i}E_{i}
=\varnothing$ and $\bigcap_{i\neq j}E_{i}\neq\varnothing$ for every $j$. The
diagonal of $D^{n}$ is then $(E_{1},\dots,E_{n})$-essential. Proof paper:
inside \pkey{thm:stable-intersection}.

\emph{Cost of the repair.} \cite[Lem.~6.11]{ZhukStrong} has a printed proof
that cites \emph{itself} three times where it means the preceding lemma, so a
self-contained rendering must supply both. Both citations were re-read verbatim
(Lemma~6.11 at 2005.00593 p.~37, Theorem~3.7 at p.~46).

\subsection{$\mathrm D6$ --- a hypothesis that can fail, and cruciality supplies it}
\label{sec:D6}

\texttt{LEMBridgeFromRelation} requires, besides subdirectness and
rectangularity of the first and last variables, that there exist
$(b_{1},\bar a,a_{n})$ and $(a_{1},\bar a,b_{n})$ in $R$ with
$(a_{1},\bar a,a_{n})\notin R$. The proof of \texttt{LEMConnectedProperties}(a)
invokes it for every constraint along a path and never discharges that
hypothesis.

\emph{It can fail.} Machine-checked witness:
$R=\{(x,z,y)\in\mathbb Z_{4}\times\mathbb Z_{2}\times\mathbb Z_{4}:
x\equiv z\equiv y\ (\mathrm{mod}\ 2)\}$ is subdirect, $\ConOne(R,1)$ is
congruence mod $2$ and coordinate $1$ is rectangular, and the third hypothesis
is unsatisfiable. The construction then yields
$\delta(x_{1},x_{2},y_{1},y_{2})=\exists z\,R(x_{1},z,y_{1})\wedge
R(x_{2},z,y_{2})$, whose projection onto the first two coordinates
\emph{equals} $\ConOne(R,1)$ rather than properly containing it --- so $\delta$
is not a bridge, clause (c) failing. The lemma is genuinely inapplicable and
genuinely necessary.

\emph{Repair.} Cruciality supplies the hypothesis. If the third hypothesis
fails at coordinates $1$ and $n$ then $R$ is the conjunction of its two
projections,
\[
  R=\{(x,\bar a,y):(x,\bar a)\in\proj_{[n-1]}(R)\ \text{and}\
     (\bar a,y)\in\proj_{[n]\setminus\{1\}}(R)\},
\]
the inclusion $\supseteq$ being exactly the failure of the hypothesis. A
constraint that is crucial in a reduction therefore satisfies the hypothesis at
every pair of its coordinates: otherwise weakening it replaces it by
constraints whose conjunction defines the same relation, so the weakened
instance has the same solution set and in particular still has no solution in
the reduction, contradicting cruciality. On the witness above,
$\proj_{1,2}(R)$ and $\proj_{2,3}(R)$ do conjoin to $R$ --- exactly the
degenerate shape the corollary rules out.

Proof paper: \pkey{lem:connected} is stated for crucial instances. This is a
notion \cite{ZhukSimplified} inherits from \texttt{LEMCrucialMeansIrreducible}
and then stops using.

\subsection{$\mathrm D7$ --- two gaps in the main induction}\label{sec:D7}

\paragraph{(i) The endgame proves \emph{connected}; (1c) demands \emph{linked}.}
Both endgames of Case~2 finish with ``\dots which means that $\mathcal I$ is
connected and satisfies (1b) if its solution set is subdirect or (1c)
otherwise''. But (1c) requires a \emph{linked} connected subinstance, and only
connectedness was established.

\emph{Repair, via irreducibility.} Suppose $\mathcal I$ is connected with
non-subdirect solution set. Take $\mathcal J=\Upsilon=\mathcal I$, legitimate
since an instance is a weakening, hence an expanded covering, of itself. It
remains to see $\mathcal I$ is linked. If it were not, $\mathcal I$ would
witness the failure of its own irreducibility: its variables are among its own,
each constraint is the projection of itself onto all its variables, it is not
fragmented (connectedness of the constraint-adjacency graph forces constraints
to share variables), it is not linked, and its solution set is not subdirect
--- precisely conditions (i)--(v) of the definition. Since $\mathcal I$ is
irreducible by standing hypothesis, it is linked. Two lines, and they explain
why irreducibility is a hypothesis of the theorem at all.

\paragraph{(ii) Case~1 derives (1c) for the wrong reduction.}
Case~1 obtains a $1$-consistent $D^{(2)}\subte{T}D^{(1)}$, shows $\mathcal I$
is crucial in $D^{(2)}$, and closes with ``applying the inductive assumption to
$D^{(2)}$ we derive the required conditions''. Conditions (1a) and (1b) mention
no reduction and transfer verbatim. (1c) does not: it asserts an expanded
covering \emph{crucial in the ambient reduction}. The inductive hypothesis
delivers $\mathcal K$ crucial in $D^{(2)}$; the goal needs a covering crucial
in $D^{(1)}$, and the containment runs the wrong way --- an instance with no
solution in the smaller $D^{(2)}$ may have one in $D^{(1)}$.

\emph{First, the induction structure, which is load-bearing and unstated.} The
proof opens ``We prove the claim by induction on the size of $D^{(1)}$.
\emph{Let us prove (2) first.}'' That ordering is not stylistic. Part~(2) at a
given measure uses (1) only at strictly smaller measure, whereas part~(1) at
that measure uses \textbf{(2) at the same measure} --- Case~1 does exactly this
at \texttt{main:3507}. So the induction is: strong induction on
$k=\mu(D^{(1)})$; at each level establish (2) for \emph{all} instances, then
(1) for all instances. That is acyclic: $(2)_{k}\leftarrow(1)_{<k}$,
$(1)_{k}\leftarrow(2)_{k},(1)_{<k}$.

\emph{Why that alone does not close the gap.} One would like to apply (2) to
$\mathcal K$ and conclude contrapositively that $\mathcal K^{(1)}$ has no
solution. But $\mathcal K$ is an expanded \emph{covering}, so it has more
variables than $\mathcal I$, and under $\mu=\sum_{x}|D^{(1)}_{x}|$ its measure
is \emph{larger}, not equal. The device that works for a weakening does not
work for a covering.

\emph{The repair, which avoids re-entering the induction.} Put
$\mathcal L=\mathcal K\wedge\mathcal I$. It is an expanded covering of
$\mathcal I$, since $\mathcal K$ is one, $\mathcal I$ is a weakening of itself
hence one, and the union of two expanded coverings is one. $\mathcal L^{(1)}$
has no solution, since $\mathcal L$ contains every constraint of $\mathcal I$
and $\mathcal I$ is crucial in $D^{(1)}$. So $\mathcal L$ may be weakened to an
instance $\mathcal L'$ crucial in $D^{(1)}$. Every constraint of $\mathcal K$
survives that weakening: each is crucial in $D^{(2)}$, so weakening it yields a
solution in $D^{(2)}\subseteq D^{(1)}$, and weakening it inside any weaker
instance yields only more solutions; so each is already crucial in $D^{(1)}$ at
every stage, and the process only weakens constraints that are not. Therefore
$\Upsilon\subseteq\mathcal K\subseteq\mathcal L'$, its properties being
intrinsic to $\Upsilon$ and untouched, and $\mathcal L'$ is an expanded
covering of $\mathcal I$ crucial in $D^{(1)}$. No appeal to the induction at a
larger measure, so the measure survives intact.

Proof paper: \pkey{thm:main-inductive} and \pkey{haz:case1c}.

\begin{entry}[The measure, and a rejected objection]\label{ent:measure}
$\mu(D)=\sum_{x}|D_{x}|$ is well founded for every appeal the induction makes.
An earlier draft of this entry, and the external review independently, held
that it is not --- on the ground that the induction is applied at expanded
coverings, where the variable set grows and $\mu$ grows with it. \textbf{That
is wrong, and the enumeration settles it.}

The proof of \texttt{THMMainInductiveCSPClaim} runs from \texttt{main:3431} to
\texttt{main:3858} and makes exactly seven appeals to its own inductive
assumption: \texttt{main:3439}, \texttt{3507}, \texttt{3511}, \texttt{3536},
\texttt{3689}, \texttt{3714}, \texttt{3724}. Each is to a \emph{weakening} of
the current instance --- so $\Var$ can only shrink --- with a reduction
strictly smaller in $\mu$: $D^{(2)}$ in the first four, and $D^{(B,\bot)}$ in
the last three, the latter being ``the maximal $1$-consistent (probably empty)
reduction \emph{for $\mathcal I$}'' below $D^{(B,\top)}$
(\texttt{main:3554}), which cuts $D^{(1)}_{z}$ to a proper
$B<^{D_{z}}_{T(\sigma)}D^{(1)}_{z}$. The three in the
$\mathcal B_{0}\neq\varnothing$ branch draw their instances from $\Omega$,
whose \emph{condition 1} (\texttt{main:3629}) is ``every instance in $\Omega$
is a weakening of $\mathcal I$''. Expanded coverings occur in that proof only
as outputs --- alternative (1c), and the auxiliary $\bot(\cdot)$,
$\Delta(\cdot)$ and $\Theta$ --- which carry no well-foundedness obligation.

What misled both readings is \texttt{main:3913}, ``By Theorem
\texttt{THMMainInductiveCSPClaim} applied to $\mathcal J'$ and
$D^{(G,\bot)}$'', where $\mathcal J'$ is a weakening of an expanded covering.
That line is at \texttt{main:3913}, inside \texttt{THMPCDoesnotKillAllSolutions}
(\texttt{main:3860--3976}), a \emph{separate} theorem proved afterwards and
consuming the main claim as a black box; the main proof does not cite it back,
so there is no mutual induction. Applying a proved theorem to a covering is
free. Proof paper: \pkey{haz:the-induction}(c).
\end{entry}

\subsection{$\mathrm D8$ --- minimal-containing-a-point versus minimal-outright}
\label{sec:D8}

At \texttt{main:3526} the main induction picks, for each variable, ``the
minimal $D^{(2)}_{x}\le_{\MT T}D^{(1)}$ containing $s(x)$'', and two lines later
invokes a lemma whose hypothesis (5) requires $D^{(2)}_{x}$ to be \emph{a
minimal $\MT T$ subuniverse} --- minimal outright. Those are not the same
condition. The lemma bridging them has been deleted, along with its citation:

\begin{quote}\footnotesize\ttfamily
3526\quad For every variable x choose the minimal D\^{}(2)\_x <=\_\{MT\} D\^{}(1) containing s(x).\\
3527\quad \%By Lemma \textbackslash ref\{LEMMinimalContainingIsMinimal\}\\
3528\quad \%\$D\^{}\{(2)\}\_\{x\}\$ is a minimal \textbackslash mathcal\{M\}T subuniverse.\\
3529\quad By Lemma \textbackslash ref\{LEMMultiTypeStillStable\}
\end{quote}

The deleted statement (\texttt{main:1864}) was: if $C$ is inclusion-minimal with
$b\in C$ and $C\le_{\MT T}B$, then no $C'\subsetneq C$ has $C'\le_{\MT T}B$.

\emph{Repair.} Fix $T\in\{\TL,\TPC,\TD\}$ and $B\lll A$, let $\Sigma$ be the
finite set of congruences that are $T$-dividing for $B$, and for $b\in B$ put
$M(b)=B\cap\bigcap_{\sigma\in\Sigma}[b]_{\sigma}$, discarding any $\sigma$ with
$B\cap[b]_{\sigma}=B$. Then $M(b)$ is the least $\MT T$ subuniverse of $B$
containing $b$; and if $c\in M(b)$ then $[c]_{\sigma}=[b]_{\sigma}$ for every
$\sigma\in\Sigma$, so $M(c)=M(b)$. Consequently an inclusion-minimal $\MT T$
set containing $b$ equals $M(b)$ and admits no proper $\MT T$ subset: any such
$C'$ is nonempty, and any $c\in C'$ has $M(c)=M(b)\subseteq C'$.

The content is the second sentence: $b\mapsto M(b)$ is a closure whose image
consists of pairwise disjoint minimal sets, because membership in $M(b)$ is
decided blockwise and blocks of an equivalence relation either coincide or are
disjoint. Proof paper: \pkey{lem:minimal-consistent} and
\pkey{haz:minimality}.

\subsection{$\mathrm D{11}$--$\mathrm D{13}$ --- three defects, one missing hypothesis}
\label{sec:D11}

\S5.4 of \cite{ZhukSimplified} is the only part of \S5 that manipulates bridges
as objects in their own right rather than consuming them, and all three of its
defects are about which auxiliary property survives a construction. All three
are cured by one hypothesis, which the source already has in its vocabulary:
\texttt{LEMPCCongruencePropertyInductiveStep} (\texttt{SS:1379}) carries, as
hypothesis (2), exactly \emph{pair-reflexivity} --- $(a,b,a,b),(b,a,b,a)\in
\delta$ for every $(a,b)\in\proj_{1,2}(\delta)$.

\paragraph{$\mathrm D{11}$: the symmetrisation formula.}
The proof of \texttt{LEMBlockOf\brk GoodBridge\brk DoesNotHaveBAC} opens by symmetrising $\delta$:
\begin{quote}\footnotesize\ttfamily
1154\quad \textbackslash sigma(x\_1,x\_2,x\_3,x\_4) = \textbackslash exists x\_5 \textbackslash exists x\_6\\
1155\quad \textbackslash delta(x\_1,x\_2,x\_5,x\_5)\\
1157\quad \textbackslash wedge \textbackslash delta(x\_3,x\_4,x\_5,x\_6).
\end{quote}
Clause (d) of the bridge definition (\texttt{main:1258}) is
``$(a_{1},a_{2},a_{3},a_{4})\in\delta$ implies
$(a_{1},a_{2})\in\sigma_{1}\iff(a_{3},a_{4})\in\sigma_{2}$''. The first
conjunct has $x_{5}$ in both of the last two places, so $(x_{5},x_{5})\in 0_{A}$
forces $x_{1}=x_{2}$. Hence $\proj_{1,2}(\sigma)\subseteq 0_{A}$, while clause
(c) demands $\proj_{1,2}(\sigma)\supsetneq\sigma_{1}=0_{A}$: \emph{the relation
defined is never a bridge}, for any $\sigma_{1}$, and every later line of the
proof is false of it.

It is a one-character slip. The intended relation is $\delta\ast\delta^{-1}$,
and the source writes exactly that, correctly, $210$ lines later at
\texttt{SS:1365--1367}, for the same purpose.

\paragraph{$\mathrm D{12}$: bridges and abelian groups, false as stated.}
\texttt{LEMNiceBridgeGivesAbelianGroup} (\texttt{SS:1103}): for $\sigma$ a
congruence and $\delta$ a bridge from $\sigma$ to $\sigma$ with
$\proj_{1,2}(\delta)=\bcol\delta=A^{2}$ and $\delta$ symmetric, there is an
abelian group $G$ with $(A/\sigma;\delta/\sigma)\cong(G;x_{1}-x_{2}=x_{3}-x_{4})$.

\emph{Counterexample.} $\mathbf A=\mathbf Z_{3}$ with
$w(x_{1},\dots,x_{4})=x_{1}+x_{2}+x_{3}+x_{4}\bmod 3$ --- idempotent since
$4\equiv 1$, and a special WNU, so $\mathbf A\in\mathcal V_{4}$. Take
$\sigma=0_{\mathbf A}$ and
$\delta=\{(x_{1},x_{2},x_{3},x_{4})\in\mathbb Z_{3}^{4}:
x_{1}-x_{2}+x_{3}-x_{4}=0\}$. Verified by exhaustion: $\delta$ is a subuniverse
of $A^{4}$; clause (d) holds; $\proj_{1,2}(\delta)=\bcol\delta=A^{2}$; $\delta$
is symmetric; and none of the six bijections $\mathbb Z_{3}\to\mathbb Z_{3}$
carries $\delta$ to $\{x_{1}-x_{2}=x_{3}-x_{4}\}$. Since every abelian group of
order $3$ is $\mathbb Z_{3}$, no $G$ works.

\emph{What is true.} $\delta$ contains all $(a,a,b,b)$, so it is the preimage
of some $K\le G\times G$ under $(x_{1},x_{2},x_{3},x_{4})\mapsto
(x_{1}-x_{2},x_{3}-x_{4})$. Clause (d) makes $K$ meet $G\times 0$ and
$0\times G$ only at the origin, and subdirectness makes it a graph, so
$\delta=\{\varphi(x_{1}-x_{2})=x_{3}-x_{4}\}$ for an automorphism $\varphi$.
Symmetry says $\varphi^{2}=\mathrm{id}$; Zhuk's conclusion is the case
$\varphi=\mathrm{id}$, and the counterexample is $\varphi=-\mathrm{id}$, which
is why it needs $p$ odd. The step that breaks is ``composing the bridge
$\delta$ with itself we get a bridge $\delta_{0}\supseteq\delta$'':
$\delta\ast\delta=\{\varphi^{2}(x_{1}-x_{2})=x_{3}-x_{4}\}$, which for
$\varphi=-\mathrm{id}$ does not contain $\delta$, and the closing count
$|\delta|=|\delta_{0}|=|A|^{3}$ needs the containment it does not have.

\paragraph{$\mathrm D{13}$: linkedness is asserted, not shown.}
\texttt{SS:1273} puts $\delta'=\delta\cap B^{4}\cap(\cover\sigma\times
\cover\sigma)$ for $B$ a block of $\operatorname{LeftLinked}(\cover\sigma)$ and
says ``Since $\bcol\delta\supseteq\cover\sigma$, $\delta'$ satisfies all the
conditions of Lemma \texttt{LEMBlockOfGoodBridgeDoesNotHaveBAC}''. Two of those
conditions are that $\proj_{1,2}(\delta')$ properly contains $0_{B}$ and is
\emph{linked}. What $\bcol\delta\supseteq\cover\sigma$ gives is the third
condition. Nothing offered bears on linkedness --- and it does not follow: the
witness $(x_{3},x_{4})$ supplied by minimality of $\cover\sigma$ lies in
$\cover\sigma$ but in no particular block, and if it lands outside $B$ for
every choice then $\proj_{1,2}(\delta')$ collapses to $0_{B}$ and $\delta'$ is
not even a bridge. The same gap recurs at \texttt{SS:1290}.

\paragraph{The single cure, and why it costs nothing.}
Pair-reflexivity kills all three. $\mathrm D{13}$: $\cover\sigma\subseteq
\proj_{1,2}(\delta)$ by minimality of $\cover\sigma$, since $\proj_{1,2}(\delta)$
is a subalgebra of $A^{2}$, stable under $\sigma$ by clause (a) and properly
containing $\sigma$ by clause (c); so pair-reflexivity puts
$(x_{1},x_{2},x_{1},x_{2})$ in $\delta'$ for every
$(x_{1},x_{2})\in\cover\sigma\cap B^{2}$, giving $\proj_{1,2}(\delta')=
\cover\sigma\cap B^{2}$ exactly --- which is linked precisely because $B$ is a
block. $\mathrm D{12}$: pair-reflexivity \emph{is} $\varphi=\mathrm{id}$, and it
also repairs the proof directly, since $\delta\subseteq\delta\ast\delta$ by
taking $(x_{1},x_{2})$ itself as the intermediate pair. $\mathrm D{11}$:
$\delta\ast\delta^{-1}$ is pair-reflexive by construction, so the corrected
formula delivers what the lemma needs.

And it costs nothing, because the two lemmas that need the hypothesis are each
invoked exactly once in the whole paper --- \texttt{LEMNontrivialReflexive\-Bridge\-Implies}
at \texttt{SS:1368}, on a $\delta\ast\delta^{-1}$, and
\texttt{LEMNiceBridgeGivesAbelianGroup} from inside that call.

Proof paper: \pkey{def:pair-reflexive}, \pkey{lem:symmetrise},
\pkey{lem:nice-bridge-abelian}, \pkey{lem:nontrivial-bridge},
\pkey{rem:phi-twist}.

\subsection{$\mathrm D{14}$ --- the box form cannot reach the subrelation}
\label{sec:D14}

\texttt{SS:1394}, inside Case~1 of
\texttt{LEMPCCongruencePropertyInductiveStep}: with $B\subt{T}A$, $|B|>1$ and
$\delta'=\delta\cap B^{4}$,

\begin{quote}\small
``By Lemma \texttt{LEMBACenterLinkedness}
$\operatorname{RightLinked}(\proj_{1,2,3}(\delta'))=B^{2}$.''
\end{quote}

\texttt{LEMBACenterLinkedness} (\cite[Prop.~2.15(i)]{BartoKozik}) restricts a
binary relation by a \emph{box} $B_{1}\times B_{2}$. Applied with
$W=\proj_{1,2,3}(\delta)$, $B_{1}=P\cap B^{2}$, $B_{2}=B$ it yields linkedness
of $\proj_{1,2,3}(\delta)\cap B^{3}$, which contains
$\proj_{1,2,3}(\delta\cap B^{4})$ and can be strictly larger: a tuple of
$\delta$ whose first three coordinates lie in $B$ need not have its fourth
there. The two really do differ --- for the meet-semilattice on $\{0,1,2\}$
with $B=\{0,1\}$, the subalgebra $K=\{(0,0),(1,2)\}$ of $A^{2}$ has
$\proj_{1}(K)\cap B=\{0,1\}$ but $\proj_{1}(K\cap B^{2})=\{0\}$. Getting
$\delta\cap B^{4}$ out of the box form instead would need linkedness of
$\delta$ in the split $(1,2)\mid(3,4)$ or $(1,2,3)\mid(4)$, and hypothesis (3)
gives neither: a path in the coarser bipartite graph does not lift, because
consecutive edges through a third coordinate may use different fourth
coordinates. The lemma's own hypothesis $R\cap(B_{1}\times B_{2})\sdle
B_{1}\times B_{2}$ is also never discharged.

\emph{The assertion is true; the citation is to the wrong form of the
principle.} What closes it is the \emph{subrelation} form: if
$W\sdle\mathbf P\times\mathbf A$ is linked and $\varnothing\neq W'\subte{T}W$
with $T\in\{\TBA,\TC\}$, then $W'$ is linked over its own projections. That
follows from tools already present in \S5 --- pp-propagation, intersection, and
factoring --- together with one elementary fact the paper never states:
$0_{\mathbf C}$ is not an absorbing subuniverse of $\mathbf C^{2}$ when
$|C|\ge2$, because absorption at position $i$ forces the term not to depend on
argument $i$, at every $i$, hence to be constant, and idempotence then forces
$|C|=1$.

The subrelation form also discharges, for free, the subdirectness hypothesis
the box form needed. Proof paper: \pkey{lem:no-absorbing-diagonal},
\pkey{lem:absorb-linked-sub}, \pkey{haz:box-vs-sub}.

\emph{Contrast.} The analogous step in \texttt{LEMPCBridgesAreTrivial} Case~2
\emph{is} justified by the source --- by first proving the closure property
$\xi\cap(A^{2}\times B\times A)\subseteq B^{4}$ (\texttt{SS:1548--1554}), which
makes $\proj_{1,2,3}(\xi\cap B^{4})=\proj_{1,2,3}(\xi)\cap(A^{2}\times B)$.
Here $B$ is a strong subuniverse rather than a linkedness block, no closure
property is available, and absorption has to do the work instead.

\subsection{$\mathrm D{15}$ --- read $\delta$ for $\sigma$, and restore stability}
\label{sec:D15}

\texttt{SS:2208--2224}. A congruence $\delta\supseteq\sigma$ is chosen maximal
with $|B/\delta|>1$, and $B/\delta$ is shown BA and centre free. To show
$\delta$ irreducible the proof supposes $\delta=S_{1}\cap\dots\cap S_{k}$ with
$S_{i}\supsetneq\delta$, finds an $i$ with $B^{2}\not\subseteq S_{i}$, and then:

\begin{quote}\small
``By Lemma \texttt{LEMLeftLinkedStayFull}
$\operatorname{LeftLinked}(S_{i}\cap(B\times B))/\sigma=(B/\sigma)^{2}$. Hence
$(S_{i}\cap(B\times B))/\sigma$ is a linked relation, which by Lemma
\texttt{LEMLinkedImpliesBACenter} implies the existence of a BA or central
subuniverse on $B/\delta$.''
\end{quote}

Two things are wrong, and they are the same thing twice. The relation is
factored by $\sigma$ while the conclusion is about $B/\delta$; those do not
match, since $(S_{i}\cap B^{2})/\sigma$ lives on $B/\sigma$. And properness is
not established: \texttt{LEMLinkedImpliesBACenter} needs a \emph{proper}
subdirect relation, $B^{2}\not\subseteq S_{i}$ gives only
$S_{i}\cap B^{2}\subsetneq B^{2}$, and factoring can close that gap.

\emph{Repair.} Read $\delta$ for $\sigma$ throughout the last four lines, which
is plainly what was meant --- $B/\delta$ was shown BA and centre free two lines
earlier for exactly this use, and the stated conclusion already says
$B/\delta$. Properness then follows, but only from a clause the proof drops
when it introduces the $S_{i}$: the definition of irreducible
(\texttt{main:1237}) requires the $S_{i}$ to be \textbf{stable under $\delta$}.
Stability makes $S_{i}$ a union of products of $\delta$-blocks, so
$(S_{i}\cap B^{2})/\delta=(B/\delta)^{2}$ would force $B^{2}\subseteq S_{i}$.
That clause is the only part of the definition of irreducibility doing work
here. Proof paper: \pkey{lem:main-existence}, \pkey{haz:main-existence}.

This propagates: \texttt{LEMUbiquity} rests on this lemma, so $\mathrm D{15}$
reaches one of the headline properties.

\subsection{$\mathrm D{16}$ --- the conclusion is the whole lemma, and is asserted}
\label{sec:D16}

\texttt{SS:2743}: $\sigma$ linear on $\mathbf A\in\Vn$ with
$\cover\sigma=A^{2}$ implies $\mathbf A/\sigma\cong\mathbf Z_{p}$. The proof
produces the perfect-linear witness $\zeta\le\mathbf A\times\mathbf A\times
\mathbf Z_{p}$ with $\proj_{1,2}(\zeta)=A^{2}$, sets $\xi(x,z)=\zeta(x,a,z)$,
and then says in full: ``Then $\xi$ is a bijective relation giving an
isomorphism $\mathbf A/\sigma\cong\mathbf Z_{p}$.'' That is the entire content
of the lemma. What $\zeta$ gives directly is only $\xi^{-1}(0)=[a]_{\sigma}$.

\emph{Repair, by a different route.} The lemma follows from a statement already
repaired here. A bridge $\delta$ from $\sigma$ to $\sigma$ with
$\bcol\delta\supsetneq\sigma$ exists by the linear criterion; replace it by
$\delta\ast\delta^{-1}$, symmetric and pair-reflexive
(\S\ref{sec:D11}). Then \texttt{LEMNontrivialReflexiveBridgeImplies}(3)
applies, and $\cover\sigma=A^{2}$ is a congruence with the single block $A$, so
$\mathbf A/\sigma\cong\mathbb Z_{p}^{m}$ as a group. Now $m=0$ gives
$\sigma=A^{2}$, impossible because an irreducible congruence is proper
(\S\ref{sec:leg-emptyfamily}); and $m\ge2$ contradicts irreducibility, because
$\mathbf A/\sigma$ is affine, so the coordinate congruences
$\theta_{i}=\{(x,y):x_{i}=y_{i}\}$ are linear subspaces of
$(\mathbb Z_{p}^{m})^{2}$, hence subuniverses of $(\mathbf A/\sigma)^{2}$, each
properly above $0_{\mathbf A/\sigma}$ with $\bigcap_{i}\theta_{i}=
0_{\mathbf A/\sigma}$. Hence $m=1$. That the group isomorphism is an
isomorphism \emph{of algebras} is the proposition of
\S\ref{sec:leg-Zp}, and without it the statement does not typecheck against
$\mathbf Z_{p}\in\Vn$.

This route uses only in-paper material, where the source's route imports
\cite[Cor.~8.17.1]{ZhukJACM} --- so the repair is \textbf{one fewer external
dependency} than the source. Proof paper: \pkey{lem:linear-on-top},
\pkey{haz:on-top}.

\begin{entry}[Salvage: the source's $\zeta$ really is a function]\label{ent:zeta}
Worth keeping, because it repairs half of the original argument and is needed
if the perfect-linear machinery is used elsewhere. Put
$t(x,y,z)=w(x,y,\dots,y,z)$ with $n-2$ copies of $y$; on $\mathbb Z_{p}$ this
is $x-y+z$, because $p\mid n-1$. Applying $t$ to three tuples of $\zeta$ over
the same $(a_{1},a_{2})$ gives, by idempotence, that the fibre is closed under
$x-y+z$, hence is a coset of a subgroup of $\mathbb Z_{p}$: empty, a singleton,
or all of $\mathbb Z_{p}$. It is not all of $\mathbb Z_{p}$, since that fibre
would contain $0$ and a nonzero element, making $(a_{1},a_{2})$ both
$\sigma$-related and not; and it is nonempty on $\cover\sigma$ since
$\proj_{1,2}(\zeta)=\cover\sigma$. So $\zeta$ is the graph of a homomorphism
$G\colon\mathbf A^{2}\to\mathbf Z_{p}$ with $G^{-1}(0)=\sigma$.
\end{entry}

\subsection{$\mathrm D{17}$ --- maximal multi-type extension. \textbf{Open}}
\label{sec:D17}

\texttt{SS:3113}. $\sigma$ is maximal with $(C_{1}\circ\sigma)\cap
B_{2}=\varnothing$; Case~1 supposes $E\subt{\TS}B_{1}/\sigma$ and argues:

\begin{quote}\small
``By Theorem \texttt{THMMainStableIntersection} $E\cap B_{2}/\sigma\neq
\varnothing$ and $E\cap C_{1}/\sigma\neq\varnothing$, hence
$C_{1}/\sigma\cap B_{2}/\sigma\neq\varnothing$, which contradicts our
assumptions.''
\end{quote}

The two nonemptinesses are fine --- though they come from the intersection
property applied to the BA-and-central $D\subseteq E$, not from the theorem
cited. \textbf{But ``hence'' does not follow}: two sets each meeting $E$ need
not meet each other. The maximality of $\sigma$, which would be the natural
extra ingredient, is not used anywhere in Case~1.

\paragraph{Why the source's Case~1 cannot be patched in place.}
The configuration \emph{reproduces itself} on $D$: one has $D\lll\bar A$,
$\bar C_{1}\cap D\lll D$, $\bar B_{2}\cap D\lll D$, still disjoint, with
$D\cap\bar B_{2}\neq\varnothing$. Every consequence of the two facts survives
the descent, so no iteration of them closes the case. The descent is not
useless --- $C_{1}\cap D\le_{\MT T}^{A}D$ properly, so the hypotheses are
inherited and an induction on $|B_{1}|$ is available --- but it shrinks the
ambient, while the conclusion's congruence family is indexed by the ambient
($\cover\omega\supseteq B_{1}^{2}$) and the single call site
(\texttt{main:2685}) consumes the original $B_{1}$ in one of its two branches.

\paragraph{The published repair, and why it fails.}
An earlier draft of this document passed to a \emph{minimal} subfamily
$F\subseteq\{\bar C_{1}^{1},\dots,\bar C_{1}^{t}\}$ with $\bigcap F\cap
\bar B_{2}=\varnothing$, ``which exists because the whole family works'', and
then used the stable-intersection theorem --- whose leave-one-out hypothesis is
exactly minimality, and whose four conclusions have no room for $\TS$ --- to
show no member of $F$ is $\TS$-typed, so that the case evaporates rather than
being contradicted.

The first line is wrong, and an external review found it. What the hypothesis
gives is $\bigl(\bigcap_{j}C_{1}^{j}\bigr)/\sigma\cap\bar B_{2}=\varnothing$,
whereas $\bigl(\bigcap_{j}C_{1}^{j}\bigr)/\sigma\subseteq
\bigcap_{j}\bigl(C_{1}^{j}/\sigma\bigr)$ with the inclusion generally strict.
So $\bigcap F$ can acquire $\sigma$-classes that meet $\bar B_{2}$, and the
subfamily need not exist. This is precisely the image-versus-intersection
failure that \S\ref{sec:leg-quotient} is about --- catalogued in this document
and then walked into. The lesson is recorded rather than tidied away.

\paragraph{Where it stands.}
\emph{Two things have changed since the repair was withdrawn.}

First, the image-versus-intersection failure is now known \emph{not} to occur
among the type-$T$ factors. If $\bar C_{1}^{j}\subt{T}\bar B_{1}$ then the
``moreover'' clause of \texttt{LEMFactorByDelta} gives
$\sigma\cap B_{1}^{2}\subseteq\sigma_{j}\cap B_{1}^{2}$, and a two-line
argument then shows
$\bigcap_{j\in F}\bar C_{1}^{j}=\bigl(\bigcap_{j\in F}C_{1}^{j}\bigr)/\sigma$
for $F$ the set of such $j$. So the trap that killed the published repair is
confined to the $\subt{\TS}$ and $=\bar B_{1}$ factors.

Second, given that, everything else in the argument runs: the proof paper
proves \pkey{lem:maximal-mult} in full from the single extra hypothesis
$(P\circ\sigma)\cap B_{2}=\varnothing$, where
$P=\bigcap_{j\in F}C_{1}^{j}$. No minimal subfamily, no $\TS$-freeness split,
and the ambient stays $\bar B_{1}$, which is what the conclusion needs.

So the whole of $\mathrm D17$ is now one sentence: \emph{the type-$T$ factors
alone already cut $B_{2}$ out}. That is strictly stronger than what the other
hypotheses give, since $C_{1}\subseteq P$.


The saturation identity $\delta\cap B^{2}\subseteq\sigma_{j}\cap B^{2}$ that
makes images commute with intersection is the ``moreover'' clause of the
factorisation lemma, and it is available for a factor precisely when that
factor lands in the type-$T$ alternative --- not when it lands in $\subt{\TS}$.
So the $\TS$ factors are the obstruction both to clause (m) of propagation and
to the minimal-subfamily argument, and the argument is circular as written.

One fragment survives, and it points at the target. A factor with
$\bar C_{1}^{j}\subt{\TS}\bar B_{1}$ contains a BA-and-central $\bar D$ of
$\bar B_{1}$, which meets $\bar B_{1}\cap\bar B_{2}$ by the strong-nonemptiness
lemma; so $\bar C_{1}^{j}\cap\bar B_{2}\neq\varnothing$ --- \emph{the $\TS$
factors never help cut $\bar B_{2}$ out}. What is missing is that the type-$T$
factors alone already cut it out. If that is proved, saturation applies to
them, the intersection of their images is the image of their intersection, and
the rest of the argument runs with the ambient still $\bar B_{1}$.

Proof paper: \pkey{lem:maximal-mult}, marked \emph{open} in its status table.
Nothing downstream of that lemma should be written until this is settled.

\subsection{$\mathrm C{10}$ --- a citation that drops a case}\label{sec:C10}

\texttt{SS:208} states:

\begin{quote}\small
``Suppose $R\sdle\mathbf A\times\mathbf B$, $C=\{c\in A\mid\forall b\in B:
(c,b)\in R\}$. Then one of the following holds: (1) $C$ is a central
subuniverse of $\mathbf A$; (2) $\mathbf B$ has a nontrivial binary absorbing
subuniverse.''
\end{quote}

\cite[Thm.~6.15]{ZhukStrong} (p.~38) is verbatim this, except for a third item:
``3. $\mathbf B$ has a nontrivial \emph{projective} subuniverse.'' The word
``projective'' does not occur anywhere in the 2024 source ---
\texttt{grep -i projective} over all four of its \TeX{} files returns nothing
--- so the notion was removed from the paper, not merely from this citation.

\emph{The two-case statement is true, and the elimination is an argument Zhuk
himself runs elsewhere without ever connecting it to Theorem~6.15.} The
standing hypothesis (\texttt{main:1123}, restated at \texttt{main:1641}
precisely for the statements \S5 proves) makes $\mathbf B$ Taylor. By
\cite[Lem.~3.4]{ZhukStrong}, a nontrivial projective subuniverse that is not
binary absorbing produces an essentially unary $\mathbf U\in\mathrm{HS}(\mathbf
B)$ of size $\ge2$ --- which a WNU forbids. Zhuk runs this same move on this
same trichotomy in the proof of his Theorem~4.15, step (4)$\Rightarrow$(5):
``In case (5) Lemma 3.4 gives us a nontrivial binary absorbing subuniverse,
which is a strong subuniverse, or an essentially unary algebra $U\in HS(A)$,
which contradicts condition (4).''

The elimination should be stated as a lemma with a proof rather than smuggled
into a citation, and the proof paper does that, with a direct two-line argument
for ``a WNU forbids an essentially unary quotient'' in place of the source's
route through WNU-blockers.

\emph{But the two-case form is true only because $\varnothing$ counts as
central.} Under our reading (\S\ref{sec:leg-empty}) it is false:
$\mathbf A=\mathbf B=\mathbf Z_{p}$ with $R$ the graph of $x\mapsto x+1$ is
subdirect, has $C=\varnothing$, and $\mathbf Z_{p}$ has no nontrivial binary
absorbing subuniverse at all. So a third alternative $C=\varnothing$ must be
added. It is harmless: each of the nine call sites exhibits an element of $C$
first.

Proof paper: \pkey{lem:no-unary-quotient}, \pkey{lem:central-relation},
\pkey{haz:central-relation}.

\begin{entry}[Two obligations at each of the eight contrapositive call sites]
\label{ent:c10-sites}
The call sites are \texttt{SS:268, 551, 606, 843, 928, 1008, 2024, 2079}. Eight
run the lemma contrapositively against a \emph{BA and center free} algebra,
defined at \texttt{main:1363} as having ``no proper nonempty binary absorbing
subuniverse or proper nonempty central subuniverse''. Two obligations are
therefore live at each and discharged at none explicitly.

\emph{Properness of the centre.} Nonempty central is not enough; $C\neq A$ must
be shown. It is available at every site, and at three of them the witness is
printed two lines earlier --- the non-containments at \texttt{SS:533, 824, 903}
--- but never connected to the appeal. Sites \texttt{606, 1008, 2079} instead
\emph{use} the improper case, concluding that $R$ is full; there the split is
explicit.

\emph{Power to base.} Alternative (2) yields a binary absorbing subuniverse of
a \emph{power}; the sites assert one on the base. That is
\texttt{LEMBACenterSOnPowerImplies}, cited at \texttt{SS:2024} and omitted at
\texttt{551, 843, 928}.
\end{entry}

\subsection{$\mathrm F$ --- a full fibre from a surjective projection}\label{sec:F}

Subsubcases 1B3 (\texttt{SS:545}), 2A3 (\texttt{SS:837}) and 2B3
(\texttt{SS:921}) each contain, verbatim modulo names:

\begin{quote}\small
``Since $\proj_{1,2,\dots,|A|}(R')=(B/\delta)^{|A|}$, there exists
$d\in B_{m-1}/\sigma_{m}$ such that $(B/\delta)^{|A|}\times\{d\}\subseteq R'$.''
\end{quote}

Surjectivity of a projection does not produce a full fibre, and as stated this
is a non sequitur.

\emph{It is nevertheless true, for a reason that explains the arity.} Each such
$R'$ is cut out by a conjunction of unary and binary conditions on its first
$|A|$ entries together with a condition relating each entry to the last
coordinate --- at \texttt{SS:837}, $\forall i,j:(a_{i},a_{j})\in\sigma\cap
\omega$, $\forall i:a_{i}\in B_{m-1}$, and $(a_{i},b)\in\sigma_{m}$ --- and any
such condition survives substituting entries for one another. And
$|B/\delta|\le|A|$ because $B\subseteq A$. So the tuple enumerating $B/\delta$
with repetition already lies in $(B/\delta)^{|A|}$, its witness $d$ serves, and
every other tuple is a coordinate-substitution instance of it.

\textbf{The exponent $|A|$ is load-bearing}: it must be at least the number of
blocks, so that one tuple can enumerate them. A rendering that normalises the
arity away breaks all three proofs. Proof paper: \pkey{lem:full-fibre},
\pkey{haz:full-fibre}.

\subsection{$\mathrm D{18}$ --- the definition of a linear congruence is ill-typed}
\label{sec:D18}

\texttt{main:1373}, clause (3) of the definition of a linear congruence, reads
verbatim:

\begin{quote}\small
``there exist prime $p$ and $S\le(\sigma^{*})^{[4]}$ such that for any block
$B$ of $\sigma^{*}$ there exists $n\ge0$ with
$(B/\sigma;S\cap(B/\sigma)^{4})\cong(\mathbb Z_{p}^{n};x_{1}-x_{2}=x_{3}-x_{4})$.''
\end{quote}

$S$ is a relation on \emph{elements} of $A$ --- it is a subalgebra of
$(\sigma^{*})^{[4]}\subseteq A^{4}$ --- while $(B/\sigma)^{4}$ is a set of
$4$-tuples of \emph{$\sigma$-classes}. The two do not intersect: they are sets
of different things, and $S\cap(B/\sigma)^{4}$ is empty on inspection, or
meaningless, depending on how literally one reads it.

\emph{Repair.} Cut down inside $A$ and then take the image:
$(S\cap B^{4})/\sigma$. That is well formed for any $S$, and it is plainly what
is meant, since the isomorphism is with a relation on $\mathbb Z_{p}^{n}\cong
B/\sigma$.

The alternative repair --- restrict the image relation $S/\sigma$ to
$(B/\sigma)^{4}$ --- gives the same relation, because $S$ is stable under
$\sigma$ in each variable and its tuples meet no other block. But that
stability is a \emph{consequence} of clause (3) (it is what makes $S$ a bridge,
the source's own following remark), so it is not available while clause (3) is
being read, and a definition may not lean on its own consequences. Proof paper:
\pkey{def:linear-congruence}, \pkey{haz:linear-typing}, and the propagation
into \pkey{lem:linear-bridge}.

Note that the corresponding clause of \texttt{LEMNontrivialReflexiveBridgeImplies}
is typed \emph{correctly} in the source --- it writes
$(B/\sigma;(\delta\cap B^{4})/\sigma)$ --- so the slip is local to the
definition.

\section{Statements restated}\label{sec:restated}

Five statements of the source are restated in the proof paper, not because they
are wrong but because our conventions make the printed form say something else
--- or because the source's own proof gives more than its statement.

\begin{center}\footnotesize
\begin{tabular}{@{}p{0.24\textwidth}p{0.32\textwidth}p{0.34\textwidth}@{}}
\hline
\textbf{Statement} & \textbf{Why} & \textbf{Restatement} \\
\hline
\texttt{LEMCentral\brk RelationImplies} (\texttt{SS:208}) &
  True in the source only because $\varnothing$ counts as central, and the
  projective alternative is dropped (\S\ref{sec:C10}) &
  ``$C=\varnothing$, or $C$ is central in $\mathbf A$, or $\mathbf B$ has a
  proper nonempty binary absorbing subuniverse'' \\
\texttt{LEMLinked\brk ImpliesBACenter} (\texttt{SS:222}) &
  ``There exists a BA or central subuniverse'' is vacuous: every algebra is a
  binary absorbing subuniverse of itself &
  Insert \emph{proper nonempty}; that is how it is used at \texttt{SS:1282} \\
The $\subt{\TS}$ clause (\texttt{main:1522}) &
  Vacuously true for $D=\varnothing$ &
  Require the witness $D$ nonempty \\
\texttt{LEMIntersectALL}(i) &
  Its proof (\texttt{SS:2356}) establishes $B\cap D\dlll^{A}D$ and then weakens
  it; three later proofs use the stronger form while citing (i) &
  State (i) as $B\cap D\dlll^{A}D$ and derive $B\cap D\dlll A$ as a corollary \\
\texttt{CORReverseHomomorphism} &
  Its type list is $\{\TBA,\TC,\TS,\TL,\TD\}$ where the stable-intersection
  theorem uses $\{\TBA,\TC,\TS,\TL,\TPC\}$ &
  Fix one convention; harmless, since $\TD$ subsumes $\TL$ and $\TPC$ \\
\hline
\end{tabular}
\end{center}
\clearpage
%% ---------------------------------------------------------------- audit2-readings.tex
\section{The legislated readings}\label{sec:readings}

Ten places where \cite{ZhukSimplified} admits more than one reading. The choice
is invisible in prose and forced in a formalization, and in six of the ten it is
not really a choice: one reading makes some statement of the source true or
some proof valid and the other does not, so the ``convention'' is a
reconstruction. Two of the ten are not conventions at all --- one is a theorem
and one is a missing lemma --- and only two are free.

\begin{center}\small
\begin{tabular}{@{}llp{0.46\textwidth}@{}}
\hline
\S & \textbf{Item} & \textbf{Status} \\
\hline
\ref{sec:leg-Zp} & the algebra $\mathbf Z_{p}$ & \textbf{a theorem}, proved \\
\ref{sec:leg-cover} & $\cover\sigma$ is a tolerance & decided by internal evidence \\
\ref{sec:leg-empty} & the empty set & decided; and it bites three statements \\
\ref{sec:leg-emptyfamily} & the empty family is allowed & decided; the two formulations agree only under it \\
\ref{sec:leg-chain} & $\lll$ carries its chain & decided; \S5.5 is the sharpest instance \\
\ref{sec:leg-dimension} & ``dimension'' & \textbf{free} \\
\ref{sec:leg-linked} & ``linked instance'' & \textbf{free} \\
\ref{sec:leg-quotient} & quotients of subsets & decided; and see $\mathrm D17$ \\
\ref{sec:leg-notice} & sentences beginning ``Notice that'' & decided; promote to lemmas \\
\ref{sec:C10} & the dropped projective alternative & \textbf{a missing lemma}, proved \\
\hline
\end{tabular}
\end{center}

\subsection{The algebra $\mathbf Z_{p}$: not a convention, a theorem}
\label{sec:leg-Zp}

\texttt{main:1115} declares \emph{``In the paper every algebra $\mathbf Z_{p}$
belongs to $\Vn$ for a fixed $n$, hence the algebra $\mathbf Z_{p}$ is uniquely
defined''}, with $w^{\mathbf Z_{p}}=x_{1}+\dots+x_{n}\bmod p$. That operation is
idempotent only when $n\equiv1\pmod p$, so the declaration is a constraint on
$p$, and the source never says where the constraint comes from.

It comes from \emph{specialness}. If $\mathbf U$ is affine over $\mathbb Z_{p}$
with $|U|>1$ and $w$ is idempotent, weak near-unanimity and special of arity
$n$, then $w$ acts as $\sum_{i}c_{i}x_{i}+d$; idempotence gives
$\sum_{i}c_{i}=1$ and $d=0$; the WNU identities force all $c_{i}$ equal to a
single $c$ with $nc=1$; and specialness compares the coefficient of $y$ on the
two sides of $w(x,\dots,x,y)=w(x,\dots,x,w(x,\dots,x,y))$ to give $c^{2}=c$, so
$c\in\{0,1\}$, and $c=0$ contradicts $nc=1$. Hence $c=1$, $n\equiv1\pmod p$,
and $w$ acts as the sum.

\emph{Rendering.} Keep the two roles of $\mathbf Z_{p}$ apart: the abelian
\emph{group} $(\mathbb Z_{p};+,-)$, which supplies the relation
$x_{1}-x_{2}=x_{3}-x_{4}$ in the definition of a linear congruence and carries
no arity constraint; and the \emph{algebra}
$\mathbf Z_{p}=(\mathbb Z_{p};x_{1}+\dots+x_{n})\in\Vn$, which does. Only the
second occurs at \texttt{main:1300} ($\zeta\le\mathbf A\times\mathbf A\times
\mathbf Z_{p}$, a subalgebra of a product, which needs a common signature) and
at \texttt{main:3484}. Proof paper: \pkey{prop:p-divides},
\pkey{cor:Zp-in-Vn}, \pkey{haz:Zp}.

The proposition is not decoration: it is needed twice, and in both places its
absence leaves a statement that does not typecheck --- it is what makes the
definition of a perfect linear congruence well formed, and what upgrades the
group isomorphism of $\mathrm D16$ to an isomorphism of algebras.

\subsection{$\cover\sigma$ is a tolerance}\label{sec:leg-cover}

$\cover\sigma$ is defined (\texttt{main:1246}) as the minimal
$\delta\le\mathbf A\times\mathbf A$ with $\delta\supsetneq\sigma$ and $\delta$
stable under $\sigma$. Nothing in that makes it transitive.

\emph{Rendering.} $\cover\sigma$ is a reflexive symmetric subalgebra of
$\mathbf A^{2}$ and \emph{not} a congruence. Three facts need one line each and
none is in the source: the family of such $\delta$ is closed under
intersection unless the intersection is $\sigma$, which is exactly what
irreducibility forbids, so a least element exists; $(\cover\sigma)^{-1}$ is
another such relation, so symmetry follows by minimality --- used silently at
\texttt{SS:1538} (``switching $x_{1}$ and $x_{2}$''); and reflexivity is
$\sigma\subseteq\cover\sigma$.

\emph{The internal evidence.} That $\cover\sigma$ is a congruence is item (2)
of the definition of a linear congruence (\texttt{main:1371}) and conclusion
(1) of \texttt{LEMNontrivialReflexiveBridgeImplies} --- i.e.\ it is a property
of linear congruences, proved, not a property of irreducible ones. Typing it as
a congruence would make item (2) vacuous and collapse the linear/PC
distinction. Proof paper: \pkey{prop:cover}, \pkey{rem:cover-not-cong},
\pkey{haz:irreducible}.

\subsection{The empty set}\label{sec:leg-empty}

Read syntactically, $\varnothing$ is a subuniverse of every idempotent algebra,
is vacuously absorbing ($t(\varnothing,\dots,A,\dots,\varnothing)\subseteq
\varnothing$ holds because the left side is empty) and vacuously central (the
clause quantifies over $a\in A\setminus C$ against $\Sg(\varnothing)=
\varnothing$).

\emph{Rendering.} Follow \texttt{main:1629} --- \emph{``we do not allow empty
subuniverses''} --- and make $\lll$, $\subt{T}$, $\subte{T}$ relations between
\emph{nonempty} sets, with dotted variants for the empty case. Three statements
must then be restated, because they are true in the source only under the
vacuous reading; they are the first three rows of \S\ref{sec:restated}.

\subsection{The empty family is allowed}\label{sec:leg-emptyfamily}

Is $\sigma=A^{2}$ irreducible? The source gives two formulations of
irreducibility (\texttt{main:1237}), and \emph{they agree only under the
empty-family reading}, which settles the question.

\begin{itemize}\tightlist
\item Formulation 1: $\sigma$ is not an intersection of \emph{other} stable
  subalgebras of $\mathbf A^{2}$. The empty intersection is $A^{2}$, so
  $\sigma=A^{2}$ is reducible.
\item Formulation 2: there are no $S_{1},\dots,S_{k}\le\mathbf A/\sigma\times
  \mathbf A/\sigma$ with $0_{\mathbf A/\sigma}=\bigcap_{i}S_{i}$ and
  $S_{i}\neq0_{\mathbf A/\sigma}$. With $k=0$ this says
  $0_{\mathbf A/\sigma}\neq(A/\sigma)^{2}$, i.e.\ $\sigma\neq A^{2}$.
\end{itemize}

Forbid the empty family and formulation~2 no longer excludes $\sigma=A^{2}$
while formulation~1 still does.

\emph{Rendering.} Allow $k=0$. Then an irreducible congruence is automatically
proper, $\cover\sigma$ exists, and \cite{ZhukJACM}'s word ``proper'' --- which
\cite{ZhukSimplified} drops --- is recovered as a consequence rather than
reinstated by hand. Proof paper: \pkey{def:irreducible},
\pkey{prop:cover}(a).

\subsection{$\lll$ carries its chain}\label{sec:leg-chain}

$C\lll^{A}B$ is defined (\texttt{main:1588}) as the \emph{existence} of a chain
$C=B_{m}\subt{T_{m}}\dots\subt{T_{1}}B_{0}=B$; several \S5 proofs then speak of
\emph{``the dividing congruences coming from $C\lll^{A}B$''}
(\texttt{main:1600}), which is a function of the chain, not of the pair.

The sharpest instance is the $\TD\times\TD$ case of
\texttt{LEMTwoStableIntersection} (\texttt{SS:1840}). The step works only if
the chain witnessing $C_{2}\lll\mathbf A$ is the chain witnessing
$B_{2}\lll\mathbf A$ extended by the single step $C_{2}\subt{\TD(\sigma_{2})}
B_{2}$ --- that is what makes the dividing congruences of $C_{2}\lll\mathbf A$
equal to those of $B_{2}\lll\mathbf A$ together with $\sigma_{2}$, which is the
whole content of the appeal. The same case also applies the clause to $C_{2}$
rather than to $B_{2}$; against $B_{2}$ it does not follow, since both
alternatives are consistent with the hypotheses and the ``moreover'' clause
then says nothing.

\emph{Rendering.} $\lll$ carries its witness: a list of pairs (subset, typed
step with its congruence). Statements that mention ``the congruences coming
from $C\lll B$'' take the chain as an argument, extension of a chain by one
step is an operation on the data, and the propositional $\lll$ is its
truncation, used only where no congruence is named. Proof paper:
\pkey{def:lll}, \pkey{haz:lll-chain}.

\subsection{``Dimension'': free}\label{sec:leg-dimension}

$\prod_{i}\mathbb Z_{q_{i}}$ with distinct primes is not a vector space over
anything, so it has no dimension. Use the composition length of the finite
abelian group, which for $\mathbb Z_{p}^{n}$ is $n$ and agrees with the source
wherever the source is meaningful, and prove additivity along short exact
sequences. Nothing in \S5 uses the notion; it is confined to the codimension-one
argument. \textbf{The legislated reading is now unnecessary there:} the repaired
Step~2 never counts a dimension. It shows $L$ is the graph of an affine map and
reads $\Delta$ off as a coset of $\ker\lambda$ for a surjection $\lambda$ onto
$\mathbf Z_{p}$, and $\lambda$ kills every factor of order coprime to $p$ for
order reasons alone. ``Codimension one'' survives only as ``$\ker\lambda$ has
index $p$''. The notion is therefore retained for the source's benefit and
consumed nowhere. Proof paper: \pkey{lem:subuniverse-coset},
\pkey{haz:step2-linear-algebra}.

\subsection{``Linked instance'': free}\label{sec:leg-linked}

Two inequivalent definitions coexist: global connectivity of the value graph,
in the informal claim, and the per-variable condition in \S3.1. The second does
not imply the first, because a fragmented instance is per-variable linked in
each component.

\emph{Rendering.} Use the \S3.1 definition and carry non-fragmentation as a
separate hypothesis wherever the informal reading was doing that work. This is
not cosmetic: the retracted recursion-depth claim (\S\ref{sec:refuted-depth})
turned on exactly this pair of conditions at the \textsc{SolveNonlinked} call
site. Proof paper: \pkey{def:instance-linked}, \pkey{haz:linked}.

\subsection{Quotients of subsets}\label{sec:leg-quotient}

$B/\sigma$ denotes the \emph{image} $\{b/\sigma:b\in B\}$ in $\mathbf A/\sigma$,
for $B$ any subset of $A$ and $\sigma$ a congruence on $\mathbf A$ --- not the
quotient of $B$ by a congruence of $\mathbf B$. Consequently
$(C_{1}\cap\dots\cap C_{t})/\delta\subseteq\bigcap_{i}C_{i}/\delta$ always,
with equality \emph{false} in general.

\emph{Rendering.} Define the image operation once, prove only the inclusion,
and flag every site that needs the reverse. The reverse inclusion is
established on the fly inside clause (fm) of Propagation --- correctly, from
the saturation identity $\delta\cap B^{2}\subseteq\sigma_{i}\cap B^{2}$ --- and
must be extracted as a standalone lemma with its hypothesis visible. Proof
paper: \pkey{haz:quotient-rel}, and the (fm) clause of
\pkey{lem:propagation}.

This is the item the invalid repair of $\mathrm D17$ walked into
(\S\ref{sec:D17}). It is worth restating what that costs: cataloguing a trap is
not the same as being immune to it, and the only defence that worked was
someone else reading the argument.

\subsection{Sentences beginning ``Notice that''}\label{sec:leg-notice}

Four items, three of them from the source's own ``Notice that'' sentences, all
used later as lemmas.

\begin{enumerate}\alphlabels\tightlist
\item \texttt{main:1378}: the relation $S$ from item (3) of the definition of a
  linear congruence is a bridge from $\sigma$ to $\sigma$ with
  $\bcol S=\proj_{1,2}(S)=\proj_{3,4}(S)=\cover\sigma$. \textbf{Checked,
  true.} Clause (4) is $x_{1}-x_{2}=0\iff x_{3}-x_{4}=0$ inside each block;
  clause (3) holds because some block has $m\ge1$, else $\cover\sigma=\sigma$;
  and $\bcol S=\cover\sigma$ because $x-x=y-y$ always. Consumed at
  \texttt{SS:1360}. Proof paper: \pkey{lem:linear-bridge}.
\item \texttt{main:1386}: $\sigma$ is irreducible, PC or linear iff
  $0_{\mathbf A/\sigma}$ is. Needed to justify the ``factorize by $\sigma$ and
  assume $\sigma=0$'' opening used in three \S5.4 proofs. Proof paper:
  \pkey{lem:irred-quotient}.
\item \texttt{main:1550}: the two formulations of $\TS$-freeness agree.
  Trivial on unfolding $\subt{\TS}$ (\texttt{main:1522}), but it must be
  written, since $\TS$-freeness is a hypothesis of clause (m) of propagation
  modulo a congruence. Proof paper: \pkey{lem:sfree-equiv}.
\item Two assertions about rectangularity belong to the same list: that the
  parallelogram property implies rectangularity (``as it follows from the
  definitions''), and that the rectangular closure exists (asserted by the
  definite article in ``\emph{the} minimal rectangular relation''). Note also
  that ``$R$ is rectangular'' unqualified --- the form used in ``rectangular
  closure'' --- is never defined in the source; it means that every variable is
  rectangular. Proof paper: \pkey{lem:rect-basics}.
\end{enumerate}

\subsection{What legislating cost}\label{sec:leg-cost}

Six of the ten (\S\S\ref{sec:leg-cover}, \ref{sec:leg-empty},
\ref{sec:leg-emptyfamily}, \ref{sec:leg-chain}, \ref{sec:leg-quotient},
\ref{sec:leg-notice}) are decided by internal evidence. Two
(\S\S\ref{sec:leg-dimension}, \ref{sec:leg-linked}) are free choices where the
source is simply ambiguous and either reading can be made to work. Two
(\S\ref{sec:leg-Zp}, \S\ref{sec:C10}) are not conventions at all and have been
proved.
\clearpage
%% ---------------------------------------------------------------- audit3-expansion.tex
\section{Steps performed silently}\label{sec:expansion}

These are not defects. They are the places where the proof paper must write out
an argument that \cite{ZhukSimplified} performs in a clause, and they are the
bulk of the work: about two dozen of them against fifteen register entries.
They are collected here so that a reader comparing the two documents can see
where the extra text comes from and satisfy themselves that it is expansion
rather than divergence.

\subsection{Coverage of \S5}\label{sec:coverage}

\S5 of \cite{ZhukSimplified} is $3\,144$ lines, $44$ statements with proofs,
$152$ proof-dependency edges. All seven subsections have now been read.

\begin{center}\small
\begin{tabular}{@{}llll@{}}
\hline
Lines & Subsection & Proofs & Defects found \\
\hline
48--70 & Additional definitions & 0 & --- \\
71--346 & Subuniverses of types $\TBA$, $\TC$, $\TS$ & 6 & none \\
347--1021 & Intersection property & 6 & none; $\mathrm F$ localised here \\
1022--1735 & Properties of PC or linear congruences & 9 & $\mathrm D11$, $\mathrm D12$, $\mathrm D13$, $\mathrm D14$ \\
1736--1927 & Types interaction & 2 & none; one missing mirror case \\
1928--2296 & Factorisation of strong subalgebras & 5 & $\mathrm D15$ \\
2297--3144 & Proof of the remaining statements & 16 & $\mathrm D16$, $\mathrm D17$ \\
\hline
\end{tabular}
\end{center}

The distribution is not an accident. \S5.4 is the only part of \S5 that
manipulates bridges as objects in their own right rather than consuming them,
and every one of its four defects is about which auxiliary property survives a
construction --- symmetrisation, restriction to a block, restriction to a
strong subuniverse.

\subsection{The four hedges}\label{sec:hedges}

A mechanical sweep for hedging phrases over the thirteen printed pages of \S5
returns four, an unusually low density: the same regex over other expositions
in this subject fires many times per page. \textbf{All four are honest.}

\begin{center}\footnotesize
\begin{tabular}{@{}lp{0.24\textwidth}p{0.50\textwidth}@{}}
\hline
Line & In & Disposition \\
\hline
\texttt{SS:146} & \texttt{LEMBACenterSImplyFactor} &
  ``For $T=\TBA$ it is straightforward'': it is one line --- if
  $t(B,A)\subseteq B$ and $t(A,B)\subseteq B$ then the same $t$ works on
  images. Only $T=\TC$ is genuinely imported. \\
\texttt{SS:177} & \texttt{LEMBACenter\brk SOnPower\brk Implies} &
  ``For $T=\TS$ just repeat the same proof \emph{word to word}'': no repetition
  is needed --- see below. \\
\texttt{SS:1560} & before \texttt{LEMNoBridge\brk Between\brk DifferentTypes} &
  ``This case can be considered in the same way as Case 2'': it is genuinely
  the mirror. $(a,b,a,b)\in\xi$ places $b$ in the block and
  $\cover\sigma\subseteq\operatorname{RightLinked}$ places $a$ and $c$;
  writing it out is mechanical. \\
\texttt{SS:2428} & \texttt{LEMFactorByDelta} region &
  ``The inclusion $\subseteq$ is obvious'': it is. \\
\hline
\end{tabular}
\end{center}

\begin{entry}[The source undersells its own citation]\label{ent:6.24}
\texttt{LEMBACenterSOnPowerImplies} says $B\subt{T}\mathbf A^{n}$ with
$T\in\{\TBA,\TC,\TS\}$ implies some $C\subt{T}\mathbf A$, and proves it by
``For $T\in\{\TBA,\TC\}$ see \cite[Lem.~6.24]{ZhukStrong}. For $T=\TS$ just
repeat the same proof word to word replacing $\TBA$ by $\TS$.'' The cited lemma
reads: \emph{``Suppose $R$ is a nontrivial strong subuniverse of $A_{1}\times
\dots\times A_{n}$ of type $T\neq\TPC$. Then there exists $i$ such that $A_{i}$
has a nontrivial subuniverse of type $T$.''} It is already general in $T$ and
already for a product of distinct algebras --- the power is a special case ---
and its proof is \emph{type-uniform}: it takes $\proj_{1}(R)$ if that is proper
and otherwise a fibre, and neither construction mentions the type, which enters
only through the facts that projections and fibres preserve $\TBA$ and preserve
$\TC$ respectively. So a single run of the proof produces one witness carrying
\emph{every} strong type $R$ has at once; if $R$ is simultaneously binary
absorbing and central --- which is what $\TS$ means --- the witness is too. The
hedge disappears rather than being discharged three times. Proof paper:
\pkey{lem:power-to-base}, \pkey{haz:power-to-base}.
\end{entry}

\subsection{\S5.3, the intersection property}\label{sec:exp-53}

\texttt{LEMIntersectionPCLinearIsGood} (\texttt{SS:623}) is a 292-line proof
with 20 distinct citations, by a wide margin the largest and most connected in
\S5. It was read in full and is \textbf{structurally sound}. Every appeal to
the inductive hypothesis lands at strictly smaller $k+\ell$: part (s) uses (s)
and (d) at $k+(\ell-1)$; subsubcase 2B2 uses (s) at $(n-1)+k$; 2B3 uses (d) at
$(n-1)+k$; and $n\le\ell$ throughout. That is worth confirming, since a
commented-out block shows an earlier draft inducted lexicographically on
$(|A/(\sigma\cap\omega)|,\,k+\ell)$ and the live proof dropped the first
component.

Seven items for the rendering. The proof paper now carries the expansion in
full; \pkey{haz:intersection-draft} keeps the four of these that a reader of
the finished proof still needs. Every consumer of the lemma was re-read against
the written-out statement afterwards --- \pkey{cor:intersection-good},
\pkey{lem:parallelogram-D}, \pkey{lem:intersect-all}(t),
\pkey{thm:stable-intersection} in its no-$\TS$, mixed and both-$\TD$ branches,
and \pkey{lem:two-stable}(2) --- and none of them breaks. Two points are worth
recording. No call site consumes minimality of the chains, which item~5 below
drops. And \pkey{lem:intersect-all}(t) is proved \emph{from} the lemma, in both
parts, so it may not be used inside it; the proof in the proof paper does not.

\begin{enumerate}\tightlist
\item \emph{Part (s): five properness splits, not one.} Throughout part (s) the
  text asserts \emph{proper} containments where
  \texttt{LEMBACenterImplyIntersection} and \texttt{LEMBACenterSImplyFactor}
  give only $\subseteq$. None is a defect in the lemma, because every equality
  branch gives the conclusion outright, but each needs writing. Branch
  $\mathcal T_{\ell}=\TD$: if
  $(G\cap C_{\ell-1})/\omega_{\ell}=C_{\ell-1}/\omega_{\ell}$ then every
  $\omega_{\ell}$-block meeting $C_{\ell-1}$ meets $G\cap C_{\ell-1}$, in
  particular the block $E$ with $C_{\ell}=C_{\ell-1}\cap E$, so
  $G\cap C_{\ell}\neq\varnothing$. Branches
  $\mathcal T_{\ell}\in\{\TBA,\TC\}$ and $\mathcal T_{\ell}=\TS$ each carry two
  more, one per containment into $B\cap C_{\ell-1}$; there the equality branch
  makes one of the two sets contain the other and the conclusion is immediate.
\item \emph{What Lemma 6.25 is doing in part (s).} Part (s) appeals to
  \texttt{LEMBACenterS\brk Possible\brk Intersections} to conclude
  nonemptiness, from a lemma whose conclusion is ``$T_{1}=T_{2}\in\{\TBA,\TC\}$'', which by itself
  permits two disjoint binary absorbing subuniverses. The step is valid only
  because $\subt{\TBA,\TC}$ is a \emph{conjunction} (\texttt{main.tex:1567}):
  one of the two sets carries both types, so the type opposite to the other
  set's is available and equality of types fails. Proof paper:
  \pkey{lem:ba-central-meets}, \pkey{haz:ba-central-meets}.
\item \emph{Properness in 2A1 and 2B1} comes from the minimal choice of $m$
  (resp.\ $n$); nonemptiness from the constant tuples
  $(b/\delta,\dots,b/\delta)$ for $b\in B$, which lie in $S_{m,0}$ since
  $B\subseteq B_{m}$ and $(b,b)\in\sigma\cap\omega$.
\item \emph{Subsubcase 2B3} writes the first alternative as
  ``$=C_{n}/\omega_{n}$'', which is stronger than the ``size 1'' the inductive
  hypothesis supplies. Correct, because $B\cap C\neq\varnothing$ and
  $C\subseteq C_{n}$ force $B_{k}\cap C_{n}\neq\varnothing$, so the single
  block is the one defining $C_{n}$; then $B_{k}\cap C_{n-1}\subseteq C_{n}$
  and $S_{k,n-1}$, $S_{k,n}$ cut $(B/\delta)^{|A|}$ identically, contradicting
  minimality of $n$.
\item \emph{Subsubcase 2A3 reads one alternative off a two-alternative lemma.}
  ``\texttt{LEMSelf\brk IntersectionPC} implies that
  $((B\circ\delta)\cap B_{m-1})/\sigma_{m}=B_{m-1}/\sigma_{m}$'' (\texttt{SS:828})
  is not an implication: that lemma offers $B_{m-1}/\sigma_{m}$ \emph{or}
  $B_{m}/\sigma_{m}$. The second is excluded, but by an argument the text does
  not give: $B_{m}=B_{m-1}\cap E$ lies inside one $\sigma_{m}$-class, so the
  second alternative makes the last coordinate of $R'$ constant, and with
  $\proj_{1,\dots,|A|}(R')=(B/\delta)^{|A|}$ that yields
  $(B/\delta)^{|A|}\times B_{m}/\sigma_{m}\subseteq R'$ --- which the minimal
  choice of $m$ forbids. Sound, but a real step.
\item \emph{The full-fibre step} at \texttt{SS:545, 837, 921} is register entry
  $\mathrm F$ (\S\ref{sec:F}).
\item \emph{Is minimality of the chains $k,\ell$ used? No.} \textbf{Settled,
  and the hypothesis is dropped.} Writing the proof out decides it in the
  negative and shows it could not have gone the other way: the inductive
  appeals in part (s), 2B2 and 2B3 are to \emph{truncations} of the chain for
  $C$, and a truncation of a minimal chain need not be minimal, so a step
  consuming minimality would break the induction rather than use it.
\item \emph{The corollary} writes $\delta=f^{-1}(\sigma)$ with $f$ undefined;
  read $f_{1}$. Its conclusion admits ``empty'' where the lemma requires
  $B\cap C\neq\varnothing$, so the empty case needs its own line at every call
  site. Proof paper: \pkey{haz:cor-intersection}.
\end{enumerate}

\texttt{LEMSelfIntersectionPC} (\texttt{SS:419}), which subsubcase 2A3
consumes, carries a proof of the same shape and three more silent steps. (i)
Subcase 1B displays ``$(E\cap B_{k-1})/\sigma_{k}=B_{k-1}/\sigma_{k-1}$''
(\texttt{SS:481}); the right-hand side is an index slip for
$B_{k-1}/\sigma_{k}$ --- $\sigma_{k-1}$ has no role in that subcase --- which
is the typo the external review flagged. (ii) The same branch concludes
$((B_{k}\circ\delta)\cap B_{k-1})/\sigma_{k}=B_{k-1}/\sigma_{k}$ from that
identity, which needs $E$ to \emph{meet} $B_{k}$; it does, because the
$\sigma_{k}$-classes the witnesses exhaust include the class of $B_{k}$, but
nothing in the text says so. (iii) Subcase 2C proves $R$ full --- a statement
about $B_{k-1}$ --- and then asserts the conclusion, which is about $B_{k}$.
The transfer works because $B_{k}=B_{k-1}\cap E_{k}$ with $E_{k}$ a
$\sigma_{k}$-class, so a $\sigma_{k}$-move inside $B_{k-1}$ cannot leave
$B_{k}$. Proof paper: \pkey{lem:self-intersection-pc},
\pkey{haz:self-intersection}.

Three smaller lemmas in the same subsection, never opened before the final
pass, are all sound. \texttt{LEMTotallySymmetricWithoutBACenter}
(\texttt{SS:357}) compresses one step: ``by the conditions of the lemma we have
$(a,\dots,a,c),(a,\dots,a,a)\in R$'' does not follow from one application of
the shift, which deletes the \emph{last} entry and so destroys $c$. What is
really happening is that symmetry plus the shift lets one pass from a tuple
with entry multiset $M$ to any $M+\{u\}-\{v\}$ with $u,v\in M$; $k-2$ such
steps delete the $b_{i}$ and duplicate $a$. This deserves its own one-line
lemma, because the same manoeuvre is needed again in
\texttt{LEMTotallySymmetricRelationForIrreducible} --- where Case~2's ``it
remains to apply the previous lemma to $R\cap(B/\sigma)^{n}$'' needs
$\proj_{1,2}(R\cap(B/\sigma)^{n})=(B/\sigma)^{2}$, not merely
$\proj_{1,2}(R)\supseteq(B/\sigma)^{2}$. Proof paper:
\pkey{lem:multiset-shift}, \pkey{haz:shift}.

\subsection{\S5.4, the bridge lemmas}\label{sec:exp-54}

Beyond $\mathrm D11$--$\mathrm D14$, \texttt{LEMBlockOfGoodBridgeDoesNotHaveBAC}
is sound in both cases, and uses five steps without stating them.

\begin{enumerate}\tightlist
\item $\bcol\delta$ is \emph{reflexive}, because
  $0_{A}\subseteq\proj_{1,2}(\delta)\subseteq\bcol\delta$, the first inclusion
  being clause (c). Everything reflexive downstream comes from this.
\item $\proj_{1,2}(\sigma\cap E^{4})=\omega\cap E^{2}$ for a subuniverse $E$,
  which is what makes Case~1's choice of $E$ the right one. It needs
  $\sigma(x_{1},x_{2},x_{1},x_{2})$ for $(x_{1},x_{2})\in\omega$ ---
  pair-reflexivity again, from the corrected symmetrisation.
\item \emph{The $E$ of Case~1 exists.} Maximality of $E$ is never used, so it
  suffices to take the block containing $a$ and $b$ of the transitive closure
  of a compatible reflexive symmetric relation on $D$ whose components are
  those of the undirected $\omega$-graph. Note that
  $(\omega\cup\omega^{-1})\cap D^{2}$, which suggests itself, is \emph{not} a
  subalgebra --- a union of two subuniverses need not be closed; use
  $(\omega\circ\omega^{-1})\cap D^{2}$, which is compatible, reflexive and
  symmetric, and has the same components because $\omega$ is reflexive.
\item The ``without loss of generality'' at \texttt{SS:1170} needs
  $\bcol\sigma$ symmetric, which holds for $\bcol\delta\circ\bcol\delta^{-1}$.
\item $\{a\}\circ\omega$ is a subuniverse, because $\{a\}$ is one by
  idempotence, which is what makes it equal to $A$ in Case~2.
\end{enumerate}

The appeal to the relational stable-intersection corollary at \texttt{SS:1214}
is \textbf{correct but not in the corollary's shape}. The three displayed boxes
$\{a\}\times B\times C\times B$, $\{a\}\times C\times C\times C$,
$\{a\}\times C\times B\times B$ are the tight boxes of three witnesses, whereas
the corollary wants $A$ in the freed coordinate. Widening each to
$\{a\}\times A\times C\times B$, $\{a\}\times C\times C\times A$,
$\{a\}\times C\times A\times B$ gives its hypothesis with $n=3$ and all
ambients $A$, legitimate because $\lll$ is reflexive; and since all three types
are $\TC$ and $n=3$, every alternative fails. The three witnesses are
$(a,b,a,b)$, $(a,a,a,a)$, $(a,a,b,b)$. Two smaller things: $C'\subte{\TC}C$ at
\texttt{SS:1226} comes out of pp-propagation as $\subte{\TC}A$ and needs
intersection to land in $C$; and the appeal there to $\proj_{1}(\xi)=A$ should
be to $\{a\}\circ\omega=A$.

\texttt{LEMNontrivialReflexiveBridgeImplies} is otherwise sound: the uniqueness
of $\cover\sigma$, the upgrade from ``every block of
$\operatorname{LeftLinked}(\cover\sigma)$ of size $>1$ satisfies
$B^{2}\subseteq\cover\sigma$'' to $\cover\sigma=
\operatorname{LeftLinked}(\cover\sigma)$, and the prime argument all check out.
Note that the contradiction is with the \emph{minimality of $\cover\sigma$},
not with irreducibility directly; the source says ``contradicts the
irreducibility of $\sigma$''. The distinction matters where $\cover\sigma$ is
data attached to a proof of irreducibility. Proof paper:
\pkey{haz:minimality-not-irred}.

\texttt{LEMPCCongruencePropertyInductiveStep} is sound outside $\mathrm D14$,
and was checked line by line. The $|A|=1$ worry does not arise: clause (c)
forces $\proj_{1,2}(\delta)\supsetneq0_{A}$, hence $|A|\ge2$. Case~2's two
applications of the ternary absorbing operation are correct and were recomputed
coordinatewise; the step needs \emph{ternary} absorption where $T=\TC$ supplies
only centrality, so ``central implies ternary absorbing''
(\texttt{main:1350}) is an unstated citation. Case~3 needs a word for the
equality case, since the absorbing-equality lemma hypothesises a proper
containment while the conclusion is immediate when it is an equality.

\texttt{LEMPCBridgesAreTrivial} is sound, with two steps to write out, both in
Case~1.

\emph{The counting step} (\texttt{SS:1478}): ``since $\proj_{1,2}(\delta)=
\proj_{3,4}(\delta)$, for any $(a,b,c,d)\in\delta$ the elements $c/\sigma$ and
$d/\sigma$ are also uniquely determined by $a/\sigma$ and $b/\sigma$''. What
the case hypothesis gives is uniqueness in the \emph{other} direction --- it is
a statement about $\delta\ast\delta^{-1}$ and the claim is about
$\delta^{-1}\ast\delta$. The bridge between them is pigeonhole, and the cited
equality of projections is what powers it: on $\cover\sigma/\sigma$, which is
finite, $\delta$ induces a bipartite relation whose left neighbourhoods are
nonempty, pairwise disjoint (that is what the case says) and cover the right
side; equal finite cardinalities force every neighbourhood to be a singleton.

\emph{The four combinations} (\texttt{SS:1534}): the proof derives
$\proj_{1,4}(\delta)=\sigma$ or $\proj_{1,3}(\delta)=\sigma$, and by the same
argument with $x_{1},x_{2}$ switched $\proj_{2,4}(\delta)=\sigma$ or
$\proj_{2,3}(\delta)=\sigma$, then says ``This completes this case''. Two of
the four combinations are the conclusion; the other two must be excluded, and
they are, because either would give $\proj_{1,2}(\delta)\subseteq\sigma$
against $\proj_{1,2}(\delta)=\cover\sigma\supsetneq\sigma$. The reverse
inclusion in the conclusion comes from stability under $\sigma$. The
``switching $x_{1}$ and $x_{2}$'' appeal needs $\cover\sigma$ symmetric.

The auxiliary relations $\zeta_{1},\zeta_{2}$ \emph{are} defined in the source,
at \texttt{SS:1486--1497}, and the proof paper reproduces those definitions
rather than reconstructing them: $\zeta_{1}$ carries eight existentially
quantified variables and four conjuncts, and the argument is a \emph{subcase
split} on whether $\zeta_{1}$ is a bridge --- not a claim that both are. An
external review was right that an earlier draft of the proof paper invoked them
without definition, and an intermediate draft supplied two short formulas that
were not the source's.

One gap at the Case~2 call: hypothesis (3) of the bridge definition for
$\xi\cap B^{4}$ requires $\cover\sigma\cap B^{2}\supsetneq\sigma\cap B^{2}$,
which is not argued. It is not needed --- if it fails, clause (d) puts both
pairs of every tuple in $\sigma$, and linkedness on a $B$ that is not a
$\sigma$-block produces $(x_{1},x_{1},x_{3},x_{3})\in\xi$ with
$(x_{1},x_{3})\notin\sigma$ directly, which is the conclusion. But the lemma
cannot be \emph{cited} on a non-bridge; the case has to be inlined. Proof
paper: inside \pkey{lem:pc-bridges-trivial}.

\subsection{\S5.5, types interaction}\label{sec:exp-55}

Both statements are correct. \texttt{LEMBridgeBetweenCongruences} omits one
direction: clause (4) is an \emph{equivalence}, and the proof establishes only
$(x_{1},x_{2})\in\sigma_{1}\Rightarrow(y_{1},y_{2})\in\sigma_{2}$. The converse
is the mirror argument, and it is where the \emph{other} half of the hypothesis
$\omega\cap\sigma_{1}=\omega\cap\sigma_{2}$ is spent. Also unstated: clause (3)
needs $\sigma_{1}\subseteq\proj_{1,2}(\delta)$ as well as the strictness
witness, and $(a,b)\notin\sigma_{2}$ follows because
$\omega\setminus\sigma_{1}=\omega\setminus\sigma_{2}$.

\texttt{LEMTwoStableIntersection} is missing one case and six steps. The
missing case is the mirror $T_{1}=\TD$, $T_{2}\in\{\TBA,\TC\}$: the hypotheses
are symmetric so it is the same argument with the indices swapped, but nothing
says so, where \texttt{SS:1805} does say ``similarly'' for the mirror of the
$\TS$ case. The six steps, all short and all discharged by
$C_{1}\cap C_{2}=\varnothing$:

\begin{enumerate}\tightlist
\item $C_{1}\cap B_{2}\subt{T_{1}}B_{1}\cap B_{2}$ needs \emph{properness}: if
  $C_{1}\cap B_{2}=B_{1}\cap B_{2}$ then $B_{1}\cap C_{2}\subseteq
  C_{1}\cap C_{2}=\varnothing$.
\item \texttt{SS:1830}: factoring returns $\subseteq$ and can collapse a proper
  containment; properness holds because $(C_{1}\cap B_{2})/\sigma_{2}=
  B_{2}/\sigma_{2}$ would make the $\sigma_{2}$-block defining $C_{2}$ meet
  $C_{1}$.
\item \texttt{SS:1817} writes the first alternative as $(B_{1}\cap B_{2})/
  \sigma_{2}=C_{2}/\sigma_{2}$ where the clause supplies only size $1$; which
  block it is follows from $B_{1}\cap C_{2}\neq\varnothing$.
\item The $\TD\times\TD$ case must apply the clause to $C_{2}$, not $B_{2}$
  (\S\ref{sec:leg-chain}).
\item The chain for $C_{2}\lll\mathbf A$ must be \emph{chosen} to extend the
  chain for $B_{2}\lll\mathbf A$ (\S\ref{sec:leg-chain}).
\item $(c_{1},c_{2})\in\omega\setminus\sigma_{1}$: membership in $\omega$
  because each $B_{i}$ lies in a single block of every dividing congruence of
  its own chain and $c_{1},c_{2}\in B_{1}\cap B_{2}$; non-membership in
  $\sigma_{1}$ because $c_{2}\in B_{1}$ and $c_{2}\,\sigma_{1}\,c_{1}\in C_{1}$
  would give $c_{2}\in C_{1}\cap C_{2}$.
\end{enumerate}

Proof paper: \pkey{lem:bridge-between-congruences}, \pkey{lem:two-stable}.

\subsection{\S5.6 and \S5.7}\label{sec:exp-5657}

Beyond $\mathrm D15$, $\mathrm D16$ and $\mathrm D17$, the remaining sixteen
proofs are sound. The steps to supply:

\begin{itemize}\tightlist
\item \texttt{LEMCenterCanBePushedIn} takes the ``full'' alternative of the
  intersection corollary without a word. ``Empty'' is excluded by
  $R\cap(B_{1}\times B_{2})\neq\varnothing$. \emph{``Size one'' needs two
  lines}: the class would have to be $[c]_{\delta}$, since $c$ has an
  $R$-partner in $B_{1}$; and then $B_{2}''=B_{2}'\cap E$ either has
  $E=[c]_{\delta}$, putting $c$ in $B_{2}''$ against the minimal choice of
  $B_{2}'$, or misses $[c]_{\delta}$ entirely, so $R\cap(B_{1}\times B_{2}'')=
  \varnothing$ against $B_{2}\subseteq B_{2}''$. Smaller: $S'=(B_{1}/\sigma)^{
  |A_{1}|}$ holds because a \emph{single} good $c$ witnesses every tuple, and
  $S''\neq\varnothing$ is where $R\cap(B_{1}\times B_{2})\neq\varnothing$ is
  spent.
\item \texttt{LEMLeftLinkedStayFull}: Case~2's restriction to $n=2^{k}$ is not
  cosmetic --- $R_{n}=(R\circ R^{-1})^{n}$ is symmetric, so
  $R_{2n}=R_{n}\circ R_{n}^{-1}$, which is what makes
  $\operatorname{LeftLinked}$ of the auxiliary relation computable in one step.
  Unstated: the auxiliary relation is subdirect over
  $(B_{1}/\sigma)\times\proj_{2}$ of itself, not over $(B_{1}/\sigma)\times
  A_{2}$; and $R_{n}\supseteq0_{A_{1}}$ by subdirectness.
\item \texttt{LEMCongruenceEitherCutOrDoNothing} leans throughout on
  $B^{2}\subseteq\omega$ --- every $B_{i}$ in the chain lies inside a single
  block of its own dividing congruence --- which is never stated and is used at
  least four times. Also unstated: $|B/\sigma|>1$, and that a path witnessing
  linkedness must contain a single step crossing $\sigma$-classes.
\item \texttt{LEMFactorByDelta}: Case~2 is a chain of six steps of which the
  source states three, and the \textbf{``moreover'' clause is never argued} ---
  it holds because Case~2 gives $\delta\cap B^{2}\subseteq\delta'\cap B^{2}=
  \sigma\cap B^{2}$ and Case~1 never produces the type-$T$ alternative. That
  clause is the saturation identity, and it is what \S\ref{sec:D17} turns on.
\item \texttt{LEMPropagation}: clause (ft) is discharged by citing the
  factorisation lemma, which is stated only for $T\in\{\TPC,\TL\}$; the
  $\{\TBA,\TC,\TS\}$ cases need clause (fs) and a sentence. And clause (fm)'s
  $t=1$ base case is where the $\TS$-freeness hypothesis is \emph{spent} --- it
  kills the factorisation lemma's middle alternative, which a multi-type cannot
  absorb. Nothing says so.
\item \texttt{CORPropagateToRelations}(e) needs two unstated steps:
  $C\circ\delta=B\circ\delta$ is excluded because $\delta\subseteq\sigma$ puts
  $C\circ\delta$ inside one $\sigma$-block while $|B/\sigma|>1$; and the
  $\subt{\TS}$ alternative is excluded because the witness also lies in one
  $\sigma$-block.
\item \texttt{LEMMultiplyByAllLinear} invokes the stable-intersection theorem
  (\texttt{SS:3057}) for its Case~1/Case~2 dichotomy. An earlier reading of
  this entry said no theorem is consumed, on the ground that walking down the
  chain and asking where the intersection with $B_{2}$ becomes empty is
  elementary. \textbf{That is right about the dichotomy and wrong about what
  Case~2 asserts.} Case~2 does not merely produce a step; it produces a step
  \emph{of type $\TD$}, and the inductive hypothesis is then applied to it,
  which needs it to be of multi-type $T$. Two things are therefore owed and
  neither is given: that the step is not of type $\TBA$, $\TC$ or $\TS$ ---
  provable for $\TS$ from \pkey{lem:intersection-good}(s), but not for the
  other two, since $\subt{\TBA,\TC}$ there is a conjunction and one absorbing
  step supplies half of it --- and, when $T\in\{\TL,\TPC\}$, that the
  irreducible congruence carried by a $\TD$ step is of type $T$ rather than the
  other one. The second is a missing hypothesis in the induction rather than a
  suppressed step. The stable-intersection appeal is presumably aimed at the
  first, but the theorem's conclusion is about a critically inconsistent family
  having equal types, and the family here is not put in that form.
  Proof paper: \pkey{haz:multiply-citation}, \pkey{cl:multiply-step}.
\item \texttt{LEMPCOnTheTopIsEasy} rests on two things it does not state: that
  an idempotent algebra is polynomially complete iff every \emph{reflexive}
  subalgebra of a power is a conjunction of equalities, and that
  $\mathbf A/\sigma$ is BA and centre free, without which the step
  $\{a_{i}\}\subt{\TPC}\mathbf A/\sigma$ is not licensed. Proof paper:
  \pkey{haz:on-top}.
\item Typographical: \texttt{LEMUbiquity}'s ``block $D$ of $B/\delta$'' should
  read ``block $D$ of $\delta$''; \texttt{LEMCenter\brk CanBePushedIn}'s $\neq$ at
  \texttt{SS:2014} should be $\not\subseteq$;
  \texttt{CORParallelogramPropertyForD}'s undefined $C$ at \texttt{SS:1000}
  should be $C_{1}$; \texttt{LEMMaximalMultExtention}'s
  $\cover{\delta_{i}}\supseteq B_{1}^{2}$ at \texttt{SS:3128} should be
  $\cover{\delta_{i}}\supseteq(B_{1}/\sigma)^{2}$; and the instance-fragmented
  definition writes $X_{1}$ twice where the second should be $X_{2}$.
\end{itemize}

\subsection{\texttt{LEMPreserveLinkdness}: a flag withdrawn}\label{sec:exp-preserve}

Flagged in an earlier survey on the grounds that the workhorse
\texttt{LEMPreserve\brk Linkdness\brk OneStepAUX} assumes both
$R\cap(B_{1}\times C_{2})\neq\varnothing$ and $R\cap(C_{1}\times B_{2})\neq
\varnothing$ while the target lemma assumes only $R\cap(B_{1}\times B_{2})\neq
\varnothing$. \textbf{Not a defect}: the flag conflates the auxiliary lemma
with \texttt{LEMPreserveLinkdnessOneStep}, which is what the four-line proof
actually invokes and which needs only $R\cap(B_{1}\times B_{2})\neq\varnothing$
and $S\cap(C_{1}\times B_{2})\neq\varnothing$.

The four-line proof is correct and unpacks to a two-phase argument the proof
paper should write out. Decompose both multi-type reductions into chains.
\emph{Phase 1}: shrink coordinate 1 along its chain with $B_{2}$ fixed; at each
step the needed $R\cap(B_{1}^{(j)}\times B_{2})\neq\varnothing$ is the previous
step's conclusion, and the needed $S\cap(B_{1}^{(j+1)}\times B_{2})\neq
\varnothing$ follows from $S\cap(C_{1}\times C_{2})\neq\varnothing$ by
monotonicity. This yields $R\cap(C_{1}\times B_{2})\neq\varnothing$.
\emph{Phase 2}: the same on the transpose $R^{-1}$, whose rectangular closure
is $S^{-1}$, shrinking coordinate 2 with $C_{1}$ fixed --- legitimate because
$C_{1}\lll A_{1}$. Proof paper: \pkey{lem:preserve-linked}.
\clearpage
%% ---------------------------------------------------------------- audit4-evidence.tex
\section{Claims examined and withdrawn}\label{sec:refuted}

Two claims were raised, survived a round of review, and turned out to be
wrong. They are the most instructive entries in this document, because both
failed in the same way and it is the way a defect list is prone to fail:
reading the formal-looking part of a paper and not the sentence before it.

\subsection{The claimed citation cycle in \S5}\label{sec:refuted-cycle}

The claim was a cycle
$\text{Lemma }9\Rightarrow7\Rightarrow86\Rightarrow85\Rightarrow
\text{Cor }22\Rightarrow\text{Thm }21\Rightarrow\text{Lemma }8\Rightarrow
\text{Lemma }9$, closed at \texttt{SS:1214}, and that it ``has to be broken
before anything in \S5.4 or \S5.7 can be formalized''.

\textbf{Refuted mechanically.} Every \verb|\begin{proof}| in \texttt{main.tex},
\texttt{StrongSubalgebras.tex} and \texttt{necessary\brk Claims.tex} was parsed,
attributed to its statement (including the \verb|\newtheorem*| duplicates the
paper uses to restate results proved later), and its \verb|\ref|s extracted.
\textbf{58 statements with proofs, 152 edges, no cycles.}

The named edge is real: \texttt{LEMBlockOfGoodBridgeDoesNotHaveBAC} does invoke
the relational stable-intersection corollary at \texttt{SS:1214}, in its
$T=\TC$ case. But it closes nothing. That corollary's transitive proof-closure
is 25 statements and does not contain
\texttt{LEMBlockOfGoodBridgeDoesNotHaveBAC}; its direct dependencies are
Propagation and the stable-intersection theorem, and the latter's closure is 15
statements, also disjoint from it.

Consequence: \S5 can be written in dependency order, and there is no hidden
simultaneous induction to untangle. The claim would not have survived ten
minutes of parsing.

\subsection{The recursion-depth bound}\label{sec:refuted-depth}

The claim was that the depth bound $|A|+|\Gamma|$ of \cite[Lem.~5.2]{ZhukJACM}
was ``not established for the \textsc{CheckTuple} $\to$ \textsc{SolveNonlinked}
path'', and that ``without the repair the algorithm is not polynomial''.
\textbf{Withdrawn.}

Lemma~5.2 bounds the depth by arguing that every recursion site except
\textsc{CheckWeakerInstance} ``reduce[s] all domains of size greater than 1''.
At the \textsc{SolveNonlinked} site this needs the instance to be neither
linked nor fragmented --- otherwise the linked-component congruence can be
trivial on some domain, that domain does not shrink, and the bound degrades
from the constant $|A|$ to $O(|A|\cdot n)$, which is fatal, since the total
work is $\mathrm{poly}(N)^{\text{depth}}$. The source does state the hypothesis
and does discharge it. It is simply not in one place:

\begin{itemize}\tightlist
\item the precondition ``an instance that is not linked and not fragmented'' is
  in the prose introducing the function (\texttt{1704.01914:892}), and
  \emph{not} in the \texttt{Input:} line of its pseudocode;
\item it is discharged at the call site (\texttt{:635}) by ``$\Theta_{0}$ is
  not fragmented because it is crucial'' --- itself a one-line argument the
  source does not give: a fragmented unsatisfiable instance has an
  unsatisfiable part, and no constraint of the \emph{other} part is then
  crucial;
\item the step from ``not fragmented'' to ``\emph{every} domain splits'' ---
  propagate the partition along a path, which exists by non-fragmentedness and
  is value-preserving by $1$-consistency --- appears once, at \texttt{:1070},
  inside the proof of a \emph{different} lemma, and is never cited here.
\end{itemize}

So the mathematics is present and the theorem is not in doubt. What is absent
is any single place where the six recursion sites of Lemma~5.2 are discharged.
The same pattern holds for the $\Gamma$-order argument, whose case (3) needs
the projections to be unchanged under weakening, which holds because the
instance is $1$-consistent when \textsc{CheckWeakerInstance} runs --- true, and
unstated.

\emph{Chasing it down did leave one genuine hole}, and closing it needed an
argument that is not in the source. \textsc{CheckTuple} calls
\textsc{SolveNonlinked} not on $\Theta_{0}$ but on $\Theta_{0}$ with domains
restricted, and restricting can only \emph{destroy} the ``not linked''
precondition.

\begin{lemma}[The \textsc{SolveNonlinked} call site]\label{lem:solvenonlinked-callsite}
Let $\Theta_{0}$ be cycle-consistent, non-fragmented and not linked, and let
$\Theta_{0}'$ be $\Theta_{0}$ with each domain restricted to a subset
$D_{x}'\subseteq D_{x}$. Then every domain of $\Theta_{0}$ of size greater than
one is strictly shrunk by the time \textsc{SolveNonlinked}$(\Theta_{0}')$
recurses into \textsc{Solve}.
\end{lemma}

\begin{proof}
Write $\tau_{x}=\mathrm{LinkedCon}(\Theta_{0},x)$, a congruence by
\cite[Lem.~6.2]{ZhukJACM}. First, $\tau_{x}$ is proper at \emph{every}
variable. Since $\Theta_{0}$ is not linked there are at least two connected
components of the graph on $\bigsqcup_{z}\{z\}\times D_{z}$; and every
component meets every variable, because from any $(x,a)$ non-fragmentedness
supplies a path to $y$ and $1$-consistency extends $a$ along it. So each
$D_{y}$ is partitioned into at least two nonempty $\tau_{y}$-blocks, each
therefore proper.

Now split on whether $\Theta_{0}'$ is linked.

\emph{Not linked.} Then \textsc{SolveNonlinked} recurses on a component, whose
trace on each $D_{y}$ is a $\tau_{y}$-block intersected with $D_{y}'$ --- a
proper subset of $D_{y}$ by the previous paragraph.

\emph{Linked.} Then for every $z$ and all $a,b\in D_{z}'$ some path of
$\Theta_{0}'$ connects them; every such path is a path of $\Theta_{0}$, so
$D_{z}'\times D_{z}'\subseteq\tau_{z}$, i.e.\ $D_{z}'$ lies inside a single
$\tau_{z}$-block. That block is proper, so $D_{z}'\subsetneq D_{z}$ --- every
domain has \emph{already} been shrunk, at the \textsc{CheckTuple} step, before
\textsc{SolveNonlinked} was called at all.
\end{proof}

\begin{entry}[Why this is the right argument]\label{ent:callsite-moral}
Two things are worth extracting.

First, the question this replaces was the wrong one. The natural worry is
whether the restriction preserves the stated precondition ``not linked'', and
one then starts asking how the minimal linear congruence $\mathrm{ConLin}$
relates to $\mathrm{LinkedCon}$ --- a relation the source never states, and
which is not needed. The dichotomy above uses no relation between them:
\emph{either} the precondition survives and the component split does the
shrinking, \emph{or} it fails and the restriction has already done it.

Second, it follows that \textsc{SolveNonlinked}'s precondition is not required
for the depth bound at all. The function is correct on a linked instance too
--- there is then one component and it recurses on the whole of it --- and the
lemma covers that case. The precondition should be read as documentation of the
intended use, not as an obligation the caller must discharge.
\end{entry}

\subsection{Declared moot: the arity in the special-WNU lemma}\label{sec:moot}

\cite[Lem.~4.7]{MarotiMcKenzie} is cited for the existence of a special
idempotent WNU of arity $n^{n!}$ in $\Clo(w)$. The concern raised was that the
period of $y\mapsto w(x,\dots,x,y)$ need not divide $n!$. The suggested
counterexample --- $n=3$ with a $5$-cycle --- does not work, since idempotence
makes $x$ a fixed point of that map and a full $5$-cycle on five elements is
then impossible; but cycles of length up to $|A|-1$ remain possible, and
$|A|-1$ need not divide $n!$ when $|A|>n$.

\emph{Moot.} The dichotomy proof consumes only ``there exist $N$ and a special
idempotent WNU of arity $N$ in $\Clo(w)$''. The precise arity matters only for
\S4 of \cite{ZhukSimplified}, which is an independent second result and out of
scope. The proof paper states the weaker form.

\subsection{A cycle in our own arrangement}\label{sec:our-cycle}

The parser that refuted the claimed cycle in the source
(\S\ref{sec:refuted-cycle}) was pointed at the proof paper, and found one
there:
\begin{center}\footnotesize
\pkey{cor:intersection-good} $\to$ \pkey{lem:propagation} $\to$
\pkey{lem:factor-by-delta} $\to$\\[2pt]
\pkey{lem:cut-or-nothing} $\to$ \pkey{lem:leftlinked-stay-full} $\to$\\[2pt]
\pkey{lem:center-pushed-in} $\to$ \pkey{cor:intersection-good}
\end{center}
Every edge was a real citation. The cycle was \emph{ours}: the source proves
the corollary from \texttt{LEMReverseHomomorphism} (\texttt{SS:953}), a
separate and much easier statement, where our draft had folded the backward
clauses into Propagation and cited the whole thing. Citing the bundle drags in
the forward clauses, which rest on \S5.6, which rests back on the corollary.

The fix was to state the backward clauses as their own lemma, which the
external review had independently asked for on different grounds --- it noted
that ``the reverse-homomorphism lemma used for (bt) is neither stated nor
cited''. The proof paper now has \pkey{lem:reverse-hom}, and its graph is
acyclic: $97$ statements, $52$ with proofs, $125$ edges.

Two things worth drawing out. First, the source was right and we were wrong,
which is the opposite of the direction this document usually runs. Second, this
is the third time the same tool has settled a question in minutes that prose
inspection had not settled at all.

\medskip

Fixing the first cycle exposed a second, and this one the review had predicted.
It wrote: ``There is also a possible dependency-cycle danger:
\texttt{block-good-bridge} invokes stable intersection; stable intersection
invokes \texttt{no-cross-bridge}; and the most natural proof of no-cross
appears to route through the linear bridge characterization and the block-good
machinery. An explicit acyclic dependency graph is needed.'' That is exactly
right, and once the two bridge-classification lemmas were written out the
parser closed the loop:
\begin{center}\footnotesize
\pkey{cor:stable-intersection} $\to$ \pkey{thm:stable-intersection} $\to$
\pkey{lem:bridge-to-pc} $\to$ \pkey{lem:no-cross-bridge} $\to$\\[2pt]
\pkey{lem:linear-equiv} $\to$ \pkey{lem:nontrivial-bridge} $\to$
\pkey{lem:block-good-bridge} $\to$ \pkey{cor:stable-intersection}
\end{center}
It is not vicious. \texttt{LEMBlockOfGoodBridgeDoesNotHaveBAC} appeals to
stable intersection at $n=3$ with all types $\TC$, and that case rests on
\cite[Lem.~6.11]{ZhukStrong} and on nothing in the bridge theory. The proof
paper therefore states that case as its own lemma, proved before the bridges,
and both consumers cite it.

The source avoids the loop by not citing: its
\texttt{LEMPCBridgesAreTrivial} contradicts PC-ness without naming
\texttt{LEMLinearEquivalentConditions}, so the edge that closes our cycle is
absent from its graph. The dependency is nevertheless there in the mathematics
--- ``$\sigma$ is PC, so no bridge from $\sigma$ to $\sigma$ has
$\bcol\delta\supsetneq\sigma$'' \emph{is} that lemma --- which is worth
recording, because it means the earlier acyclicity verdict on the source
(\S\ref{sec:refuted-cycle}) is a statement about its citations and not about
its logical dependencies. The claim refuted there was a specific named cycle,
and that refutation stands; the stronger reading, that no cycle of any kind is
latent in \S5, was never established and is now known to be delicate.

\section{The machine checks}\label{sec:machine-checks}

Everything reducible to a finite check has had one. Each entry gives the scope,
so that a reader can see what it does and does not cover. Two of the original
claims of this document died to exactly such checks.

\begin{center}\footnotesize
\begin{tabular}{@{}p{0.20\textwidth}p{0.30\textwidth}p{0.38\textwidth}@{}}
\hline
\textbf{Check} & \textbf{Scope} & \textbf{Result} \\
\hline
Proof-dependency graph, the source & 58 statements, 152 edges over the three
  \TeX{} files & Acyclic; \S\ref{sec:refuted-cycle} \\
Proof-dependency graph, the proof paper & 97 statements, 52 with proofs, 125
  edges & Acyclic --- after one real cycle was found and fixed,
  \S\ref{sec:our-cycle}. Run by the build gate \\
Statement audit & all 81 labelled statements & 9 flagged, 8 cleared by hand;
  $\mathrm D2$ is the sole survivor \\
$\mathrm C1$: special idempotent affine WNU on $\mathbb Z_{p}$ &
  $p\in\{2,3,5,7\}$, $3\le n<12$ & exists exactly when $p\mid n-1$, and is then
  the sum \\
$\mathrm D6$: the third hypothesis can fail &
  $\mathbb Z_{4}\times\mathbb Z_{2}\times\mathbb Z_{4}$ & witness found;
  \S\ref{sec:D6} \\
$\mathrm D{12}$: counterexample & all six bijections $\mathbb Z_{3}\to
  \mathbb Z_{3}$, exhaustive over $\mathbb Z_{3}^{4}$ & none carries $\delta$
  to the difference relation \\
$\mathrm D{14}$: the subrelation form &
  nine small Taylor algebras (min-semilattices, the dual discriminator,
  $\mathbb Z_{2}$ minority, $\mathbb Z_{3}$, $\mathbb Z_{2}^{2}$, a semilattice
  square), $527+99\,477$ pairs $(W,W')$ & no counterexample; and $0_{C}$
  absorbs $C^{2}$ in no algebra tested \\
\hline
\end{tabular}
\end{center}

\begin{entry}[What the searches are and are not]\label{ent:search-limits}
A search over nine small algebras is evidence that a statement is not
\emph{obviously} false, and nothing more; the two positive entries above
($\mathrm C1$ and $\mathrm D14$) are backed by written proofs in the proof
paper, and the searches were run first, to decide whether the proofs were worth
attempting. The two \emph{negative} entries are different in kind: a
counterexample found by exhaustion over a finite structure is a proof, and
$\mathrm D6$ and $\mathrm D12$ rest on nothing else. Those two are stated in
the proof paper in a form a reader can check by hand in a few lines, which is
the right standard for a paper; the exhaustive verification is a check on the
transcription, not the argument.

Where the proof paper still says ``verified by exhaustion'' for a claim with no
short symbolic proof, that is a gap, and it is listed in
\S\ref{sec:risk}.
\end{entry}

\section{Disposition of the external review}\label{sec:review}

An external reviewer read a single-file build of an earlier draft. Its
findings, and what was done with each.

\subsection{Accepted}\label{sec:review-accepted}

\begin{itemize}\tightlist
\item \textbf{$\mathrm D17$ is broken at its first line.} Correct, and it
  reverses a claim made here. \S\ref{sec:D17}; the item is reopened.
\item The symmetrisation lemma cited bridge composition, which carries an
  irreducibility hypothesis it does not have. Fixed by proving the collapse
  identity directly.
\item The nice-bridge lemma asserted that $\delta\ast\delta$ is an equivalence
  relation, which needs a stabilised power or the Maltsev route.
\item The block-good-bridge proof treated $(\omega\cup\omega^{-1})\cap D^{2}$ as
  a subalgebra, which a union is not; use $(\omega\circ\omega^{-1})\cap D^{2}$
  (\S\ref{sec:exp-54}).
\item ``Which we may take reflexive'' is used twice with no lemma behind it.
  Accepted as a gap in our exposition, but the gap is not where either of us
  thought: there is nothing to normalise. Every consumer hypothesises
  $\bcol\delta\supsetneq\sigma$, which contains $\sigma\subseteq\bcol\delta$,
  and $(a,a)\in\sigma\subseteq\bcol\delta$ \emph{is} $(a,a,a,a)\in\delta$. So
  the bridge is reflexive already, and the proof paper says so in one line
  (\pkey{lem:bridge-reflexive}). A genuine normalisation of the same
  \emph{shape} does occur elsewhere --- ``we may assume $\bcol\delta$
  rectangular, as otherwise compose it with itself many times'', in
  \texttt{LEMNoBridge\brk Between\brk DifferentTypes} --- and that one is a
  genuine replacement of $\delta$ rather than an observation. It is now
  \pkey{lem:rectangularise}: the composite's collapse is
  $\LeftLinked(\bcol\delta)\circ\bcol\delta$, a union of full boxes, and it
  disturbs neither linking equivalence. The same proof spends rectangularity
  exactly once, and the projection step beside it named the wrong
  projection --- $\proj_{1}(\bcol\delta)$ for $\proj_{2}(\bcol\delta)$ --- which
  is what made it look like a second gap.
\item The abelian-bridge characterisation needs $|A|>1$, and ``composing enough
  times'' needs an explicit stabilised power.
\item The reason given for reflexivity in \pkey{lem:perfect-from-linked} is
  false: $(a,a)$ lying in the first projection gives some $\delta(a,a,c,d)$,
  not $\delta(a,a,a,a)$. The repair is the reviewer's: clause (d) gives
  $c\,\sigma\,d$ and stability replaces $d$ by $c$.
\item The document had the wrong shape, mixing a proof paper with a project
  status report. That is what this three-way split answers.
\item A duplicate label, several overfull boxes, a \verb|\ref| printed inside a
  \verb|\cite|, and an unidentifiable bibliography entry. All fixed; the build
  gate now fails on each class.
\end{itemize}

Several further findings concern statements the proof paper had not yet written
in full --- the induction measure across coverings (Entry~\ref{ent:measure}),
weakening's termination order, the common binary absorption witness
(\S\ref{sec:D1}), the mixed-prime linear algebra, and the definition of
$\textsc{Solve}$. They are scheduled work rather than audit entries, and they
are listed in the proof paper's status table.

\subsection{Accepted, with a correction}\label{sec:review-corrected}

The reviewer wrote that $\LeftLinked$ and $\RightLinked$ ``are used throughout
the strong-subalgebra proofs \dots but no definition is given in the
manuscript, and Brady's notes do not use those exact names''. Both halves are
true of \emph{our} draft and of \cite{BradyNotes}. But
\cite{ZhukSimplified} \textbf{does} define them, at \texttt{SS:52--66}, and in
the general $n$-ary form:

\begin{quote}\small
``For a relation $R\le\mathbf A_{1}\times\dots\mathbf A_{n}$ by
$\LeftLinked(R)$ we denote the minimal equivalence relation on $\proj_{1}(R)$
such that $(a_{1},a_{2},\dots,a_{n}),(b_{1},a_{2},\dots,a_{n})\in R$ implies
$(a_{1},b_{1})\in\LeftLinked(R)$. Similarly, $\RightLinked(R)$ is the minimal
equivalence relation on $\proj_{n}(R)$ such that
$(a_{1},\dots,a_{n-1},a_{n}),(a_{1},\dots,a_{n-1},b_{n})\in R$ implies
$(a_{n},b_{n})\in\RightLinked(R)$.''
\end{quote}

So this is not a defect of the source but an omission of ours: we reproduced
the uses without the definition. It is now
\pkey{def:linked} and \pkey{lem:linked-basics}, with the equivalence to
connectedness of the bipartite graph proved for subdirect binary relations.

The $n$-ary form settles a question that matters elsewhere. $\LeftLinked$
splits $R$ at its \emph{first} coordinate and $\RightLinked$ at its
\emph{last}, with all intermediate coordinates going to the other side --- so
$\RightLinked(\proj_{1,2,3}(\delta))$ concerns the split $(1,2)\mid(3)$ and is
a relation on the third coordinate. That is exactly the reading $\mathrm D14$
(\S\ref{sec:D14}) turns on, and it is now confirmed against the definition
rather than inferred from the uses.

\subsection{Rejected, with reasons}\label{sec:review-rejected}

\begin{itemize}\tightlist
\item \textbf{The witness in \texttt{LEMBridgeBetweenCongruences} is correct as
  printed.} The reviewer proposed reading $z_{i}=y_{i}=x_{i}$ for
  $z_{i}=y_{i}=x_{1}$. The printed choice works: it gives
  $(x_{1},x_{2},x_{1},x_{1})\in\delta$ using symmetry of $\sigma_{1}$ and
  reflexivity of $\omega$, which is what is needed for
  $\sigma_{1}\subseteq\proj_{1,2}(\delta)$. The proposed choice needs
  $(x_{1},x_{2})\in\omega$, which is not available.
\item \textbf{$\subt{\TBA,\TC}$ is a conjunction, not a disjunction.}
  \texttt{main:1567} reads ``we write $C<_{\TBA,\TC}B$ meaning that
  $C<_{\TBA}B$ and $C<_{\TC}B$''. The reviewer is right that the notation was
  never defined in our document, and it now is; but the reading is ``both''.
\item \textbf{The induction measure does survive expanded coverings, because
  the induction is not applied at one.} The reviewer's \S3, and an earlier
  draft of ours, hold that $\mu(D)=\sum_{x}|D_{x}|$ is not well founded
  because a covering has more variables. The premise --- that the induction
  appeals at a covering --- is false: all seven appeals inside
  \texttt{main:3431--3858} are to weakenings, and the covering application is
  at \texttt{main:3913}, inside the \emph{next} theorem. Entry~\ref{ent:measure}
  gives the enumeration. This is the one review finding whose rejection
  required reading the source rather than re-reading our own text.
\end{itemize}

\section{Risk register}\label{sec:risk}

What is least well tested, stated so that a later reader knows where to push.

\begin{description}\tightlist
\item[Open mathematics.] $\mathrm D17$ (\S\ref{sec:D17}) and the step
  \pkey{cl:multiply-step} inside \pkey{lem:multiply-\brk all-linear}.
  These two gate the statements that depend on them. Three items listed here in earlier drafts
  have since come off: the reflexivisation, which needs no argument at all; the
  rectangularity normalisation, which is \pkey{lem:rectangularise}; and the
  induction measure across expanded coverings (Entry~\ref{ent:measure}), which
  was never in difficulty.

\item[Least-tested repairs.] $\mathrm D7$(ii)'s reweakening and $\mathrm D6$'s
  cruciality corollary each touch the CSP-instance machinery rather than the
  algebra, and each has been checked against only the call site that motivated
  it.

\item[Least-tested positive verdicts.] The self-intersection lemma was read only
  where the central-relation call sites forced it open. \S5.3's three small
  lemmas were opened last, and one of them is what splits the main case of the
  largest proof in \S5.

\item[Unread.] \texttt{main.tex}'s CSP half --- instances, coverings,
  cruciality, the main induction, the algorithm --- has been read only where
  $\mathrm D6$, $\mathrm D7$ and $\mathrm D8$ touch it. That is roughly two
  thousand lines to which the standard of \S\ref{sec:standard} has not been
  applied, and the register above should not be read as covering them. \S4 of
  the source (\texttt{XYSymmetric.tex}, 1\,640 lines) is out of scope by
  decision.

\item[Claims resting on exhaustion alone.] None of the register's positive
  claims, but the proof paper still carries ``verified by exhaustion'' where a
  short symbolic argument would serve --- for instance the $\mathbb Z_{3}$
  counterexample of $\mathrm D12$, where setting
  $(x_{1},x_{2},x_{3},x_{4})=(0,b,b,0)$ forces $2(f(0)-f(b))=0$, so $f$ is
  constant for odd $p$ and no bijection can work. Replace them.
\end{description}

\begin{entry}[The failure mode this document is prone to]\label{ent:meta}
Three times now a claim here has been published, reviewed, and turned out to be
wrong: the citation cycle, the recursion-depth bound, and the repair of
$\mathrm D17$. In each case the error survived internal review and died to
someone reading the argument from outside. A list like this one is exactly as
reliable as the reading behind each entry, and the failure mode is not
carelessness --- it is confidence in a step that was never read at the standard
of \S\ref{sec:standard}. The entries most at risk are the ones in the
\emph{unread} row above.
\end{entry}

\clearpage
%% ---------------------------------------------------------------- bibliography.tex
% Shared bibliography for the three documents.  Kept in one file so that a
% citation key means the same thing in the proof paper, the audit and the
% blueprint.
\begin{thebibliography}{99}

\bibitem{BJK} A.~Bulatov, P.~Jeavons, A.~Krokhin,
  \emph{Classifying the complexity of constraints using finite algebras},
  SIAM J. Comput.\ \textbf{34} (2005), 720--742.

\bibitem{BradyNotes} Z.~Brady, \emph{Notes on CSPs and polymorphisms},
  arXiv:2210.07383. Cited by section title and numbered statement; the
  sections used here are ``Cores and Idempotent Reducts'' and \S3.11,
  \S3.15.

\bibitem{BartoKozik} L.~Barto, M.~Kozik, \emph{Absorbing subalgebras, cyclic
  terms, and the constraint satisfaction problem}, Log. Methods Comput. Sci.
  \textbf{8} (2012), 1:07.

\bibitem{DecidingAbsorption} L.~Barto, A.~Kazda, \emph{Deciding absorption},
  Internat. J. Algebra Comput. \textbf{26} (2016), 1033--1060.

\bibitem{Bergman} C.~Bergman, \emph{Universal algebra: fundamentals and
  selected topics}, CRC Press, 2011.

\bibitem{Bulatov} A.~Bulatov, \emph{A dichotomy theorem for nonuniform CSPs},
  arXiv:1703.03021; FOCS 2017.

\bibitem{FederVardi} T.~Feder, M.~Y.~Vardi, \emph{The computational structure of
  monotone monadic SNP and constraint satisfaction}, SIAM J. Comput.\
  \textbf{28} (1999), 57--104.

\bibitem{HobbyMcKenzie} D.~Hobby, R.~McKenzie, \emph{The structure of finite
  algebras}, Contemp. Math.~76, Amer. Math. Soc., 1988.

\bibitem{IstingerKaiser} M.~Istinger, H.~K.~Kaiser, \emph{A characterization of
  polynomially complete algebras}, J. Algebra \textbf{56} (1979), 103--110.

\bibitem{MarotiMcKenzie} M.~Mar\'oti, R.~McKenzie, \emph{Existence theorems for
  weakly symmetric operations}, Algebra Universalis \textbf{59} (2008),
  463--489.

\bibitem{MinimalTaylor} L.~Barto, Z.~Brady, A.~Bulatov, M.~Kozik, D.~Zhuk,
  \emph{Minimal Taylor algebras as a common framework for the three algebraic
  approaches to the CSP}, LICS 2021.

\bibitem{RossSlides} R.~Willard, \emph{Similarity, critical relations, and
  Zhuk's bridges}, slides, AAA98 --- 98th Workshop on General Algebra,
  Dresden, 2019.
  \url{http://www.math.uwaterloo.ca/~rdwillar/documents/Slides/willard-aaa98-handout.pdf}

\bibitem{ZhukJACM} D.~Zhuk, \emph{A proof of the CSP dichotomy conjecture},
  J. ACM \textbf{67} (2020), no.~5, 1--78; arXiv:1704.01914.

\bibitem{ZhukSimplified} D.~Zhuk, \emph{A simplified proof of the CSP dichotomy
  conjecture and XY-symmetric operations}, arXiv:2404.01080v2.

\bibitem{ZhukStrong} D.~Zhuk, \emph{Strong subalgebras and the constraint
  satisfaction problem}, J. Mult.-Valued Logic Soft Comput.\ \textbf{36} (2021);
  arXiv:2005.00593.

\end{thebibliography}

\end{document}
